\DeclareUnicodeCharacter{0301}{\'e}
\documentclass[aps,pra,reprint,groupedaddress,floatfix]{revtex4-2}
%\documentclass[aps,rmp,preprint,superscriptaddress]{revtex4-2}
%\documentclass[aps,rmp,reprint,groupedaddress]{revtex4-2}

%\usepackage[margin=1in]{geometry}
\usepackage{multirow}
\usepackage{amssymb,amsmath,amsthm,amsfonts}
\usepackage{bm}
\usepackage{textgreek}
\usepackage{mathtools}
\usepackage{enumitem}
%\usepackage[numbers,comma,sort&compress]{natbib}
%\usepackage{authblk}
\usepackage{graphicx}
\usepackage{subcaption}
\usepackage{tablefootnote} 
\usepackage[ruled,vlined,linesnumbered]{algorithm2e}
\usepackage{algorithmic}
\renewcommand{\algorithmicrequire}{\textbf{Input: }}  
\renewcommand{\algorithmicensure}{\textbf{Output: }}
\usepackage{physics}


\usepackage{footnote}
\usepackage{xcolor}
\usepackage{mathrsfs}
\usepackage{bbm}
\usepackage{makecell}
\usepackage{pbox}
\usepackage[colorlinks]{hyperref}
\hypersetup{citecolor=blue}
% \urlstyle{same}
%\captionsetup[figure]{labelfont=bf}

\newtheorem{fact}{Fact}[section] 
\newtheorem{theorem}{Theorem}
\newtheorem{definition}{Definition}
\newtheorem{lemma}{Lemma}
\newtheorem{proposition}{Proposition}
\newtheorem{example}{Example}
\newtheorem{corollary}{Corollary}
\newtheorem{remark}{Remark}
\newtheorem{assumption}{Assumption}
\newcommand{\eqn}[1]{(\ref{eqn:#1})}
\newcommand{\eq}[1]{(\ref{eq:#1})}
\renewcommand{\prb}[1]{(\ref{prb:#1})}
\newcommand{\thm}[1]{\hyperref[thm:#1]{Theorem~\ref*{thm:#1}}}
\newcommand{\cor}[1]{\hyperref[cor:#1]{Corollary~\ref*{cor:#1}}}
\newcommand{\defn}[1]{\hyperref[defn:#1]{Definition~\ref*{defn:#1}}}
\newcommand{\lem}[1]{\hyperref[lem:#1]{Lemma~\ref*{lem:#1}}}
\newcommand{\prop}[1]{\hyperref[prop:#1]{Proposition~\ref*{prop:#1}}}
\newcommand{\assum}[1]{\hyperref[assum:#1]{Assumption~\ref*{assum:#1}}}
\newcommand{\fig}[1]{\hyperref[fig:#1]{Figure~\ref*{fig:#1}}}
\newcommand{\tab}[1]{\hyperref[tab:#1]{Table~\ref*{tab:#1}}}
\newcommand{\algo}[1]{\hyperref[algo:#1]{Algorithm~\ref*{algo:#1}}}
\renewcommand{\sec}[1]{\hyperref[sec:#1]{Section~\ref*{sec:#1}}}
\newcommand{\append}[1]{\hyperref[append:#1]{Appendix~\ref*{append:#1}}}
\newcommand{\fac}[1]{\hyperref[fac:#1]{Fact~\ref*{fac:#1}}}
\newcommand{\lin}[1]{\hyperref[lin:#1]{Line~\ref*{lin:#1}}}
\newcommand{\itm}[1]{\ref{itm:#1}}
\newcommand{\bdelta}{\mathbf{\Delta}}
\renewcommand{\arraystretch}{1.25}
\newcommand{\algname}[1]{\textup{\texttt{#1}}}
\renewcommand{\i}{\mathrm{i}}
\renewcommand{\d}{\mathrm{d}}
\renewcommand{\L}{\mathcal{L}}
\newcommand{\ad}{\mathrm{ad}}
\newcommand{\T}{\mathcal{T}}
\newcommand{\B}{\Big}
\renewcommand{\P}{\mathcal{P}}
\newcommand{\simu}{\mathrm{sim}}
\newcommand{\Hc}{H_{\mathrm{control}}}
\newcommand{\He}{H_{\mathrm{eff}}}
\newcommand{\A}{\mathcal{A}}
\renewcommand{\O}{\mathcal{O}}
\newcommand{\tot}{\mathrm{tot}}
\newcommand{\G}{\mathcal{G}}
\newcommand{\E}{\mathcal{E}}
\renewcommand{\comm}{\mathrm{comm}}
\newcommand{\supp}{\mathrm{supp}}
\newcommand{\eps}{\varepsilon}
\newcommand{\mc}{\mathcal}
\newcommand{\mb}{\mathbb}
\def\>{\rangle}
\def\<{\langle}
\newcommand{\trans}{^{\top}}

\newcommand{\xz}[1]{{\color{cyan} XZ: #1}}
\newcommand{\tyl}[1]{{\color{blue} Tongyang: #1}}
\newcommand{\syz}[1]{{\color{teal} (Shengyu: #1)}}
\newcommand{\syzn}[1]{{\color{teal} #1}}
\newcommand{\sz}[1]{{\color{blue} Shuo: #1}}

\allowdisplaybreaks


\begin{document}

% Override caption formatting to force justified (not centered) captions,
% unaffected by \centering in figure environments.
\makeatletter
\long\def\@makecaption#1#2{%
  \par
  \vskip\abovecaptionskip
  \begingroup
   \small\rmfamily
   \samepage
   \leftskip\z@\rightskip\z@\parfillskip\@flushglue
   \let\footnote\@footnotemark@gobble
   \@make@capt@title{#1}{#2}\par
  \endgroup
  \vskip\belowcaptionskip
}
\makeatother

\title{Jointly Mitigating Physical and Trotter Errors in Hamiltonian Simulation with Commutator Bounds}

\author{(experiment) + theory folks}
\date{}

\begin{abstract}
Product-formula Hamiltonian simulation is attractive for near-term quantum devices because of its simple circuit structure and commutator-sensitive error scaling, but its performance is limited by both algorithmic Trotter error and physical noise. Existing mitigation strategies typically handle these two errors separately, requiring extrapolation over both the Trotter step size and the noise strength. We develop a one-dimensional joint mitigation protocol that couples the tunable noise strength to the Trotter step size as $\lambda(s)=c(sT)^{p+1}$ for a $p$-th order product formula, so that the leading physical-noise and Trotter contributions appear at the same order and can be canceled by a single Richardson extrapolation. The analysis uses a finite-order BCH truncation bound controlled by doubly right-nested commutators, avoiding the need for convergence of the full BCH series and yielding commutator-sensitive guarantees for noisy product formulas. We instantiate the bounds for $k$-local Hamiltonians and second-quantized plane-wave electronic-structure Hamiltonians, and show that under sufficiently mild local Lindbladian noise the physical noise can be mitigated together with the Trotter error with only a constant asymptotic overhead compared with noiseless Trotter extrapolation. Numerical simulations further demonstrate that the proposed one-dimensional procedure can achieve accuracy comparable to a two-dimensional Richardson baseline while using fewer extrapolation points.
\end{abstract}

\maketitle

Summary of the results: 
\begin{enumerate}
    %\item Provide an upper bound for the truncation error of the BCH formula of $\prod_v e^{\mc L_v}$ (each $\mc L_v$ is a Lindbladian) governed by the doubly right-nested commutators of $\mc L_v$.
    \item Provide a general condition involving the doubly right-nested commutators of Hamiltonians, where the Trotter error mitigation algorithm \cite{watson2025exponentially} achieves commutator and logarithm precision scaling, and show that $k$-local Hamiltonians and the second-quantized plane-wave electronic structure Hamiltonians satisfy the condition.
    %\item Propose a Lindbladian simulation algorithm with commutator and logarithm precision scaling. 
    \item Show that under mild physical noise, the Trotter error mitigation algorithm can simultaneously mitigate the physical noise incurring a constant complexity overhead. \textcolor{red}{Yicong: In addition to the number of interpolation nodes, try to prove that the overall sample complexity of ZNE procedure via 1D extrapolation is smaller than that by 2D extrapolation.} \sz{2D complexity to be added, even though maybe mild improvement.}
\end{enumerate}
All of the algorithms use at most two ancilla qubits. \sz{Confused: I doubt if this is the cost of lindbladian simulation, Trotter needs no ancilla.} 

\section{Introduction}
Hamiltonian simulation is one of the most promising candidates to demonstrate practical quantum advantage on near-term devices~\cite{kim2023evidence}, with applications ranging from condensed matter physics~\cite{childs2018toward}, quantum chemistry~\cite{mcardle2020quantum}, to high-energy physics~\cite{bauer2023quantum}.

The first explicit quantum simulation algorithm was proposed by Lloyd~\cite{doi:10.1126/science.273.5278.1073} using product formulas~\cite{suzuki1985decomposition,suzuki1991general}, and then numerous advanced techniques have been extensively developed, including linear combination of unitaries (LCU)~\cite{childs2012hamiltonian, dominic_w_berry_simulating_2015}, quantum signal processing (QSP)~\cite{Low_2017, low2019hamiltonian}, etc. Despite the inferior cost dependence in time $t$ and precision $\epsilon$, the product formula enjoys commutator scaling and does not require ancillary qubits or complex controlled gates, enabling its convenient implementation on near-term quantum devices.

\citet{watson2025exponentially} exponentially reduced the error dependence of Trotter using extrapolation. However, providing fully rigorous commutator bounds for general Hamiltonians is challenging. A key technical difficulty is that the convergence of the BCH formula, which the analysis relies on, is not guaranteed in all regimes (e.g., for general $k$-local Hamiltonians), making the derivation of commutator bounds difficult. Important progress was recently made by \citet{mizuta2025commutator}, who provided a rigorous commutator bound for the specific case of lattice Hamiltonians by using a subsystem Trotterization technique.
Furthermore, our recent work~\cite{wang2026lindbladiansimulationcommutatorbounds} derives an upper bound of the truncation error of the BCH formula for general Lindbladians in terms of the doubly right-nested commutators, which naturally provides a rigorous error bound for Trotter extrapolation for general Hamiltonians, including lattice and second-quantized Hamiltonians even with dissipation.

Besides the Trotter error, physical noise is another central obstruction on near-term devices. Quantum error correction provides a scalable route to fault tolerance~\cite{Terhal_2015}, but it requires resources far beyond what is currently available in near-term experiments. Quantum error mitigation (QEM)~\cite{cai2023quantum}, including zero-noise extrapolation (ZNE)~\cite{li2017efficient, PhysRevLett.119.180509}, probabilistic error cancellation~\cite{EndoBenjaminLi2018, berg_probabilistic_2023}, and measurement error mitigation~\cite{Bravyi2021measurement, Nation2021scalable}, is more immediately compatible with noisy devices, though it can increase the number of circuit samples needed for a fixed accuracy.


Most existing approaches treat algorithmic Trotter error and physical noise separately, for example by combining ZNE with an extrapolation in the Trotter step size or conversely, performing a two-dimensional extrapolation~\cite{vazquez2023well}. Recent work showed that the two errors can instead be coupled and mitigated by a one-dimensional extrapolation~\cite{hakkaku2025dataefficienterrormitigationphysical}. However, the existing analysis focuses on first-order Trotter formulas, uses an approximate global depolarizing-noise model, and does not provide commutator-sensitive bounds. This motivates the central question addressed in this work:
\begin{align*}
    \textit{Can we jointly mitigate physical and Trotter errors while maintaining commutator scaling?}
\end{align*}

\textcolor{red}{Tongyang: to be extended}
The main contribution of this work is to answer this question by applying the Lindbladian BCH truncation analysis of Ref.~\cite{wang2026lindbladiansimulationcommutatorbounds} to noisy product formulas. By coupling the noise strength $\lambda$ to the Trotter step size $s$ as $\lambda = c(sT)^{p+1}$ for a $p$-th order product formula, the leading physical-noise and Trotter contributions occur at the same order in $s$ and can be removed by a single Richardson extrapolation. We prove commutator-sensitive bounds for this one-dimensional joint mitigation scheme, including applications to $k$-local lattice Hamiltonians and second-quantized plane-wave electronic-structure Hamiltonians. Under sufficiently mild intrinsic noise, the physical noise can be mitigated together with the Trotter error with only a constant asymptotic overhead.


\section{Preliminaries}
% \textbf{Roadmap.}
 In this section, we introduce the notations, key concepts, and auxiliary lemmas used throughout the paper. \sz{Truncation error of BCH may become an individual section instead of a subsection.}


\subsection{Notations}
We denote the set of all positive integers by $\mb{Z}_+$. For a sequence of operators $\A_1, \ldots, \A_{\Gamma}$, we denote $\prod_{\gamma=1}^\Gamma \A_{\gamma}= \A_{\Gamma}\A_{\Gamma-1} \ldots \A_1$ and $\prod_{\gamma=\Gamma}^1 \A_{\gamma}= \A_{1}\A_{2} \ldots \A_\Gamma$. We denote the right-nested commutator $[X_1, [X_2, [\cdots[X_{n-1}, X_n]\cdots]]]$ by $[X_1, X_2, \ldots, X_n]$. We denote $\{1, 2, \ldots, n\}$ by $[n]$ and the collection of all size-$m$ subsets of $[n]$ by $\binom{[n]}{m}$. We use $\supp(X)$ to denote the set of qubits that $X$ nontrivially acts on.  For a superoperator $\mc L$, we denote the diamond norm of $\mc L$ by $\|\mc L\|$.
Define the time-ordering operator $\mc T\{\cdot\}$ as \begin{align*}
    \mc T\big\{[O_1(t_1),O_2(t_2),\ldots, O_q(t_q) ]\big\} = [O_{\pi(1)}(t_{\pi(1)}),O_{\pi(2)}(t_{\pi(2)}),\ldots, O_{\pi(q)}(t_{\pi(q)})],
\end{align*}
where $\pi\colon [q]\to[q]$ is a permutation of $[q]$ such that \[
    t_{\pi(1)}\ge t_{\pi(2)}\ge \cdots \ge t_{\pi(q)}.
    \]
For any permutation $\sigma\colon [q]\to [q]$, define the reordering operator $\mc R_{\sigma}\{\cdot\}$ as \begin{align*}
    \mc R_{\sigma}\big\{[O_1(t_1),O_2(t_2),\ldots, O_q(t_q)]\big\} = [O_{\sigma(1)}(t_{\sigma(1)}),O_{\sigma(2)}(t_{\sigma(2)}),\ldots, O_{\sigma(q)}(t_{\sigma(q)})].
\end{align*}
Throughout this paper, we denote the Lindblad dissipator with a jump operator $L$ as $\mathcal{D}[L](\rho):= L\rho L^\dagger - \frac{1}{2}\{L^\dagger L, \rho\}$.

\subsection{Quantum error mitigation}
\textit{Quantum error mitigation} (QEM)~\cite{cai2023quantum} refers to algorithmic schemes that reduce the noise-induced bias in the observable expectation value by post-processing outputs from an ensemble of circuit runs. QEM has enabled the quantum hardware in the noisy intermediate-scale quantum era to achieve practical applications \sz{may tweak: demonstrate practical quantum advantages}.

\subsubsection{Zero-noise extrapolation}
One of the most common QEM techniques is \textit{zero-noise extrapolation} (ZNE)~\cite{li2017efficient, PhysRevLett.119.180509}, which tries to obtain the zero-noise value using data points at different circuit fault rates. Specifically, we use $\rho_\lambda$ to denote the output quantum state at a given circuit fault rate $\lambda$, and model the observable average $\operatorname{Tr}[O\rho_\lambda]$ using an ansatz function $f(\lambda,\vec{\theta})$ with a set of parameters $\vec{\theta}$. Starting from the smallest circuit fault rate $\lambda_1$ that we can achieve, we boost the fault rate to obtain a set of data points $\{\lambda_m,\operatorname{Tr}[O\rho_{\lambda_m}]\}$. We then fit the ansatz to the data points to choose a set of optimal parameters $\vec{\theta^*}$. The zero-noise observable average $\operatorname{Tr}[O\rho_0]$ is estimated by
\begin{align*}
    \operatorname{Tr}[O\rho_0]\approx f(0,\vec{\theta^*}).
\end{align*}

If $\lambda$ is \textit{small}, Ref.~\cite{PhysRevLett.119.180509} showed that $\operatorname{Tr}[O\rho_\lambda]$ can be approximated using a polynomial ansatz in the same spirit as truncated Taylor series:
\begin{equation*}
\operatorname{Tr}\left[O \rho_\lambda\right] \approx f(\lambda ; \vec{\theta})=\sum_{\ell=0}^{M-1} \theta_{\ell} \frac{\lambda^{\ell}}{\ell!},
\end{equation*}
where the $(M-1)$-degree polynomial is determined by $M$ different free parameters and $f(0,\vec{\theta^*})={\theta_0^*}$. In general, the number of data points $\lambda_m$ is greater than or equal to $M$. For the simplest case where $M=2$~\cite{li2017efficient}, we can perform linear regression. When the number of data points reaches its minimum needed amount $M$, we can perform Richardson extrapolation as discussed in \sec{Richardson}. \sz{In practice, when using product formulas, we can only consider odd or even polynomials as ansatzes.}

A straightforward approach to boost the fault rate $\lambda$ is \textit{unitary folding}, which replaces a quantum circuit or gate $U$ by $U(U^\dagger U)^n$~\cite{giurgica-tiron_digital_2020, DumitrescuEtAl2018, HeEtAl2020}. This increases the noise rate by a factor of $1+2n$ without knowing the noise model. However, the circuit depth will also increase accordingly, exceeding the favorable regime for near-term devices. Alternatively, if we can completely characterize the noise model a priori, we can probabilistically insert gates to simulate the faults directly~\cite {li2017efficient}. On the other hand, knowing a concrete noise model enables us to adopt more precise ansatzes. For instance, one can rigorously prove that amplifying global depolarizing noise via unitary folding leads to the scaling form $\operatorname{Tr}[O\rho_\lambda]=a+bp^\lambda$~\cite{giurgica-tiron_digital_2020}.

% \sz{Polynomial ansatz, exponential ansatz, Richardson extrapolation. The order of the bias corresponds to the number of nodes.}

\subsubsection{Richardson extrapolation}\label{sec:Richardson}
\sz{Add intro, or remove subsestion as inline lemma: Richardson extrapolation.}
\begin{lemma}[{see \cite[Lemma A.9]{wang2025randomized}}]
\label{lem:richardson-extrapolation}
    Suppose that function $f\colon \mb R \to \mathbb{R}$ can be decomposed as 
    \begin{align*}
        f(x) = P_m(x)+R_m(x),
    \end{align*} 
    where $P_m(x)$ is a degree-$(m-1)$ polynomial and $R_m(x)$ is the remainder. Given any $t > 0$ and $s\in (0,1)$, we can choose 
    \begin{align*}
        K = \max\Big\{\frac{m}{\pi}, \frac{2t}{s}\Big\}, \quad r_j =\left\lceil \frac{K}{\sin^2(\pi(2j-1)/8m)}\right\rceil,
    \end{align*}
    and
    \begin{align*}
        s_j = \frac{t}{r_j}, \quad b_j=\prod_{\ell \neq j} \frac{1}{1-r_\ell/r_j}
    \end{align*}
    so that $1\le \|\mathbf{b}\|_1 \le C\log m$ for some absolute constant $C$, $r_j=O(\max\{m^3, m^2t/s\}/j^2)$, and $\max_{j}\{ s_j \}\le s$. Then define the $m$-term Richardson extrapolation \begin{align}
        \label{eq_richardson extrapolation}
        F^{(m)}(s)=\sum_{j=1}^m b_j f(s_j),
    \end{align}
    which satisfies \begin{align}
    \label{eq:error-F}
        |F^{(m)}(s)-f(0)|\le \|\mathbf{b}\|_1 \max_j|R_{m}(s_j)|.
    \end{align}
\end{lemma}

%\subsection{Magnus expansion}
\sz{This section might be removed or moved to appendix.}
\subsection{Technical lemmas}
\sz{Maybe move to subsubsection below.}

Let $\mc A(\tau)$ and $\mc B(\tau)$ be two continuous operator-valued functions. We rely on the following lemma to decompose the time-ordered exponential of the sum $\mc A(\tau) + \mc B(\tau)$. This result is essentially the interaction picture representation of the evolution.
\begin{lemma}[{\cite[Page 21]{dollard1984product}}]
\label{lem:interaction-pic}
    Let $\mc H(\tau) = \mc A(\tau) + \mc B(\tau)$ be an operator-valued function defined for $\tau \in \mathbb{R}$ with continuous summands $\mc A(\tau)$ and $\mc B(\tau)$. Then
    \begin{align*}
    \exp_{\mc T}\left(\int_{0}^{t}  \mc H(\tau)\, \d \tau\right) =  \begin{multlined}[t]
        \exp_{\mc T}\left(\int_{0}^{t}  \mc A(\tau)\, \d \tau\right) \\
     \cdot \exp_{\mc T}\left(\int_{0}^{t}\exp_{\mc T}^{-1}\left(\int_{0}^{\tau_1}  \mc A(\tau_2)\, \d \tau_2\right)\mc  B(\tau_1) \exp_{\mc T}\left(\int_{0}^{\tau_1}\mc  A(\tau_2)\, \d \tau_2\right)\, \d \tau_1\right).
    \end{multlined}
    \end{align*}
\end{lemma} 
We use the following lemma to bound the difference between powers of a superoperator and a quantum channel. 
\begin{lemma}[{\cite[Lemma~3]{wang2026lindbladiansimulationcommutatorbounds}}]
\label{lem:power-bound}
    Let $\mc N_1, \mc N_2$ be two superoperators, where $\mc N_1$ is a quantum channel. For any positive integer $k$, the difference between $\mc N_1^k$ and $\mc N_2^{k}$ can be bounded as
    \begin{align*}
        \|\mc N_1^k-\mc N_2^k\| \le k \max\{1, \|\mc N_2\|\}^{k-1} \|\mc N_1-\mc N_2\|.
    \end{align*}
\end{lemma}

\subsection{Truncation error of the BCH formula}
\label{sec:bch-truncate}
Let $\mc K_1, \ldots, \mc K_M$ be general Lindbladians. The BCH formula can then be written as
\begin{align}
\label{eq:def-bch}
 	e^{\mc K_M}\cdots e^{\mc K_2}e^{\mc K_1} = \exp(\sum_{q = 1}^{\infty} \Phi_q),
\end{align}
where each $\Phi_q$ is a sum of the weighted right-nested commutators of $\mc K_j$ (see \citet{arnal2021note}):
\begin{align}
\label{eq:bch}
 	\Phi_q = \sum_{\substack{p_v \ge 0,\\p_1+\cdots+p_M = q}}\frac{1}{p_1!\cdots p_M!} \sum_{\sigma\in S_q} c_{\sigma, q}\mc R_{\sigma}\{[\underbrace{\mc K_M, \ldots, \mc K_M}_{p_M}, \ldots, \underbrace{\mc K_1, \ldots, \mc K_1}_{p_1}]\},
\end{align}
where $S_q$ denotes the set of permutations of $[q]$, \begin{align}
\label{eq:def-c-sigma}
    c_{\sigma, q}=\frac{1}{q^2}\frac{(-1)^{d_{\sigma}}}{\binom{q-1}{d_\sigma}} 
\end{align} satisfying $|c_{\sigma, q}|\le 1/q^2$, and $d_{\sigma}$ denotes the number of $i\in[q-1]$ such that $\sigma(i)>\sigma(i+1)$.

Define the sum of the $q$-fold right-nested commutator norms as
\begin{align}
\label{eq:def-comm}
 	\alpha_{\rm comm}^{(q)} =\sum_{v_1, \ldots, v_q = 1}^{M} \left\|[\mc K_{v_1}, \ldots, \mc K_{v_q}]\right\|.
\end{align}
\citet{aftab2024multi} showed that the norm of $\Phi_q$ can be bounded as
\begin{align*}
 	\|\Phi_q\|\le \frac{1}{q^2} 	\alpha_{\rm comm}^{(q)},
\end{align*}
hence, the series expansion in Eq.~\eq{bch} converges if there exist two constants $J,C\ge 0$ such that
\begin{align}
\label{eq:converg-bch}
 	\sup_{q\ge J}\alpha_{\rm comm}^{(q)} \le C.
\end{align} 
A trivial bound for $\alpha_{\rm comm}^{(q)}$ is
\begin{align}
\label{eq:trivial-alpha}
 	\alpha_{\rm comm}^{(q)} \le \left(2\sum_{v = 1}^M \|\mc K_v\|\right)^q,
\end{align}
by $\|[A,B]\|\le 2\|A\|\|B\|$ and the triangle inequality. The quantity $\alpha_{\rm comm}^{(q+1)}$ also appears in the complexity of $q$-th order Trotterization \cite{childs2021theory}, where each $\mc K_v$ is a Hermitian operator in the decomposition of the simulated Hamiltonian. For some physical Hamiltonians, such as $k$-local Hamiltonians on a lattice of size $N$, where each $\mc K_v$ can be decomposed as a sum of terms acting nontrivially on at most $k$ sites, $\alpha_{\rm comm}^{(q)}$ admits a sharper bound than Eq.~\eq{trivial-alpha}. This leads to the commutator scaling of the $q$-th order Trotterization. Specifically, suppose that the maximum interaction strength of $\sum_{v = 1}^M \mc K_v$ on each site is $g$. Eq.~\eq{trivial-alpha} yields $(2Ng)^q$, while the bound considering the commutator structure of $\mc K_v$ gives $Nq!(2kg)^q = O(N(qkg)^q)$. The latter bound gives better complexity scaling for local Hamiltonians where $k\ll N$ and $q$ is a constant. 

However, in the error analysis of some advanced Trotter methods with higher precision \cite{aftab2024multi, watson2025exponentially}, bounds on $\alpha_{\rm comm}^{(q)}$ for a single $q$ do not suffice. Instead, we require the BCH formula to be convergent. The convergence condition in Eq.~\eq{converg-bch} requires $\alpha_{\rm comm}^{(q)}$ to be bounded for arbitrarily large $q$, and is thus too strict to utilize the commutator structure of $\mc K_v$, since the only valid choice to satisfy the condition is $g = 0$. \xz{add connection to trotter, and comparision to previous analysis} Concretely, in Trotter error mitigation~\cite{watson2025exponentially}, one expands the extrapolation function $f(s)$ as a power series around $s=0$, whose validity relies on the convergence of the underlying BCH series. For a $k$-local system, each generator satisfies $\mc K_v = O(s)$, yielding $\alpha_{\rm comm}^{(q)} = O(s^q q! (kg)^q)$. The factorial growth dominates the $s^q$ suppression, so the BCH series has zero radius of convergence: for any $s>0$, the series diverges. Consequently, the power expansion of $f(s)$ is merely a formal asymptotic series rather than a convergent one. This contradicts the implicit analyticity assumption in the asymptotic analysis and is precisely the gap that makes the prior analyses~\cite{aftab2024multi, watson2025exponentially} lack full rigor for general $k$-local systems. Our work~\cite{wang2026lindbladiansimulationcommutatorbounds} resolves this issue via the truncated BCH formula (\thm{general-bch-truncation}), which avoids the convergence requirement by working directly with a finite-order truncation and bounding the remainder, yielding rigorous error bounds even outside the convergence radius. The situation is the same if each $\mc K_v$ is a $k$-local Lindbladian, as the parameter $\alpha_{\rm comm}^{(q)}$ admits similar bounds (see \cor{commu-bound-sum}). 

% To bypass the divergence issue of the BCH formula in the analysis of the multiproduct formula, \citet{mizuta2025commutator} showed that for $k$-local Hamiltonians on a lattice, the truncated BCH formula could approximate the exponential product even outside its convergence radius. Consequently, they analyze the truncated BCH formula instead of the full series, thereby avoiding the convergence requirement. Their proof relies on a subsystem Trotterization technique specific to lattices, which is difficult to extend to general Hamiltonians. \sz{Ref.~\cite{wang2026lindbladiansimulationcommutatorbounds} established the BCH truncation error for general cases by constructing an equivalent continuous-time evolution with Magnus expansion.} To analyze the BCH truncation error for general cases, we construct an equivalent continuous-time evolution via the Magnus expansion
To bypass the divergence issue of the BCH formula, \citet{mizuta2025commutator} showed that, for $k$-local Hamiltonians on a lattice, a truncated BCH formula can approximate the exponential product even outside the convergence radius of the full BCH series. Their argument avoids the convergence requirement by working directly with a finite truncation, but it relies on a subsystem Trotterization technique that is specific to lattice Hamiltonians. For the general setting considered here, we instead use the truncation theorem of Ref.~\cite{wang2026lindbladiansimulationcommutatorbounds}, which is proved by constructing an equivalent continuous-time evolution and applying the Magnus expansion. This approach reduces the truncation problem to bounding the following doubly right-nested commutators.

Define the sum of doubly right-nested commutator norms as
\begin{align}
\label{eq:def-beta-bar-general}
    \bar{\alpha}_{\comm}^{(q_1, \ldots, q_{d})}
    =~ 
    \sum_{v_{1,1}, \ldots, v_{d, q_{d}} = 1}^M \bigg\| \Big[ 
         [\mc K_{v_{1,1}}, \ldots, \mc K_{v_{1,q_1}}], \ldots, 
        [\mc K_{v_{d,1}}, \ldots, \mc K_{v_{d,q_{d}}}]
     \Big] \bigg\|.
\end{align}
% For $k$-local Hamiltonians and second-quantized plane-wave electronic structure Hamiltonians, these doubly right-nested commutators admit a better scaling than right-nested ones, allowing us to derive a meaningful bound. \xz{maybe explain what is the better scaling of the doubly right-nested commutator here for some specific Hamiltonian.}
These quantities are tailored to the finite-order BCH truncation error. For $k$-local systems, the doubly nested structure can be bounded by products of local factors $(2kg)^{q_r}$ and support-growth factors $P_{r+1}$, rather than by a global norm bound such as Eq.~\eq{trivial-alpha}. For second-quantized plane-wave electronic-structure Hamiltonians, an analogous layer-counting argument yields bounds of order $\prod_r P_{r+1}q_r!n^{q_r}$. These factorized estimates are strong enough to make the truncated BCH error small in regimes where the full BCH series is not guaranteed to converge.

% \sz{Maybe it suffices to only reserve Theorem 1, removing all the former lemmas.}
We first state the general truncation theorem, and then estimate its commutator parameter for the concrete Hamiltonian classes of interest.
In the following theorem, we recall the BCH truncation bound that allows us to control the product-formula error without assuming convergence of the full BCH series.
\begin{theorem}[{\cite[Theorem~6]{wang2026lindbladiansimulationcommutatorbounds}}, BCH Truncation Bound]
 \label{thm:general-bch-truncation}
 For any positive integer $q_0$, define 
 \begin{align}
\label{eq:def-beta-comm-sum}
    \bar{\alpha}_{\comm,q_0} := \bar{\alpha}_{\comm}^{(q_0+1)} + \sum_{r = 0}^{\infty}\frac{1}{(r+1)!} \sum_{\substack{1\le p_1, \ldots, p_{r+1}\le q_0\\p_1+\cdots+p_{r+1} \ge q_0+2}}\bar{\alpha}_{\comm}^{(p_1, \ldots, p_{r+1})}.
\end{align}
 If $\bar{\alpha}_{\comm,q_0} \le 1$, the truncation error of the $q_0$-th order BCH formula is bounded as
 \begin{align*}
 \left\|\exp\Big(\sum_{q=1}^{q_0} \Phi_q\Big) - \prod_{v = 1}^Me^{\mc K_v} \right\| \le e\bar{\alpha}_{\comm,q_0}.
 \end{align*}
\end{theorem}


%%%%%%%%%%%%%%%%%%%%%%%%%%%%%%%%%%%%%%%%%%%%%%%%%%%%%%%%%%%%%%%%%%%%%%%%%%%%%%%%%%%%%%%%%%%%%%%%%%%%%%%%%%%%%
\sz{Consider move this applications to Sec 3. }
\subsubsection{Application to $k$-local systems}
\label{sec:k-local-example}

We now demonstrate the power of \thm{general-bch-truncation} by applying it to the $k$-local systems that motivated our analysis. Suppose that each generator $\mc K_v$ is a \emph{$k$-local superoperator} on a lattice $\Lambda = [N]$. Formally, this means
\begin{align*}
    \mc K_v = \sum_{\gamma = 1}^{\Gamma} \mc K_{\gamma}^{(v)}, \text{ with } |\supp(\mc K_{\gamma}^{(v)} )|\le k, \quad \forall v\in [M], \gamma\in [\Gamma].
\end{align*} 
Note that while our motivation comes from Lindbladians, the following derivation relies solely on the locality structure and holds for general superoperators.
For any site $j\in \Lambda$, we further suppose that 
\begin{align*}
    \sum_{v=1}^M\sum_{\gamma\in [\Gamma]: j\in \supp(\mc K_{\gamma}^{(v)})}\|\mc K_\gamma^{(v)}\|\le g,
\end{align*}
which is called \emph{$g$-extensiveness} of $\mc K:= \sum_{v=1}^M \mc K_v$.

A wide class of physically relevant models, specifically Lindbladians generated by sparse Pauli strings, naturally fits into this framework. In the following lemma, we translate the Hamiltonian and jump operator parameters into the locality and extensiveness of the resulting superoperators.

\begin{lemma}
\label{lem:lindblad-mapping}
    Consider a Lindbladian $\mathcal{L} = \mathcal{H} + \sum_{v=1}^M \mathcal{D}[D_v]$ acting on $N$ qubits, where
    \begin{enumerate}
        \item the Hamiltonian $H = \sum_{\gamma} h_\gamma P_\gamma$ consists of Pauli strings with weight at most $k$; and 
        \item each jump operator $D_v = \sum_{\eta} d_{v,\eta} Q_{v,\eta}$ consists of Pauli strings with weight at most $k$.
    \end{enumerate}
    Then the total Lindbladian $\mathcal{L}$ can be decomposed as a sum of \emph{$2k$-local superoperators}:
    \begin{align*}
        \mathcal{L} = \sum_{\mu} \mathcal{K}_\mu, \quad \text{with } |\supp(\mathcal{K}_\mu)| \le 2k.
    \end{align*}
    Furthermore, let the Hamiltonian extensiveness be $g_H = \max_j \sum_{\gamma: j \in \supp(P_\gamma)} |h_\gamma|$ and the dissipative extensiveness be $g_D = \max_j \sum_{v} \sum_{\eta, \lambda: j \in \supp(Q_{v,\eta}) \cup \supp(Q_{v,\lambda})} |d_{v,\eta} d_{v,\lambda}^*|$.
    The system satisfies the \emph{$G$-extensiveness} condition with parameter
    \begin{align*}
        G = 2 g_H + 2 g_D.
    \end{align*}
\end{lemma}
\begin{proof}
    The Hamiltonian part $\mathcal{H} = -i \sum_\gamma h_\gamma [P_\gamma, \cdot]$ is a sum of terms with locality $|\supp(P_\gamma)| \le k$. Since $\|[P_\gamma, \cdot]\| \le 2\|P_\gamma\| = 2$, each term contributes at most $2|h_\gamma|$ to the local extensiveness.
    
    For the dissipative part, we expand $\mathcal{D}[D_v](\rho) = \sum_{\eta,\lambda} d_{v,\eta}d_{v,\lambda}^* \mathcal{S}_{\eta,\lambda}(\rho)$, where the basis superoperator is $\mathcal{S}_{\eta,\lambda}(\rho) = Q_{v,\eta}\rho Q_{v,\lambda}^\dagger - \frac{1}{2}\{Q_{v,\lambda}^\dagger Q_{v,\eta}, \rho\}$. 
    Each $\mathcal{S}_{\eta,\lambda}$ acts non-trivially only on $\supp(Q_{v,\eta}) \cup \supp(Q_{v,\lambda})$, so its locality is bounded by $2k$.
    Regarding the norm, since $Q_{v,\eta}$ and $Q_{v,\lambda}$ are Pauli strings with unit operator norm, standard norm inequalities for unitary conjugation and anti-commutators imply $\|\mathcal{S}_{\eta,\lambda}\|\le 1 + \frac{1}{2}(2) = 2$.
    Summing the coefficients weighted by this norm bound directly yields the extensiveness parameter $G = 2g_H + 2g_D$.
\end{proof}

This lemma ensures that the results we derive for general local superoperators are directly applicable to sparse Pauli-Lindbladians by simply substituting the locality $k \to 2k$ and extensiveness $g \to G$.

For general $k$-local Hamiltonians, the standard commutator bound defined in Eq.~\eq{def-comm} is $\alpha_{\rm comm}^{(q)} = O( q! (kg)^q)$ \cite{childs2021theory,mizuta2025commutator}, which causes the BCH series to diverge \cite{mizuta2025commutator}. This scaling of $\alpha_{\rm comm}^{(q)}$ remains the same for $k$-local superoperators, as the key requirement of the proof is locality rather than Hermiticity. \citet{mizuta2025commutator} provided an upper bound on the sum of nested commutators for a single Hamiltonian. However, the proof relies essentially on the commutativity structure determined by the supports of the operators, rather than the specific forms. Therefore, this result admits a straightforward generalization to a sequence of distinct local superoperators.

\begin{lemma}[Generalization of {\cite[Lemma~7]{mizuta2025commutator}}]
\label{lem:commu-bound}
    Let $\mc K_1, \ldots, \mc K_q$ be a sequence of superoperators on a size-$N$ lattice $\Lambda$ such that $\mc K_j = \sum_{v=1}^{M} \mc K_{j,v}$ is $k$-local and $g_j$-extensive for all $j\in [q  ]$. Then the following inequality on the right-nested commutators holds: \begin{align*}
        \sum_{v_1, \ldots, v_q=1}^{M} \big\|[\mc K_{1,v_1}, \ldots, \mc K_{q',v_{q'}}, \mc A, \mc K_{q'+1,v_{q'+1}}, \ldots, \mc K_{q,v_q}]\big\|\le q!(2k)^q\Big(\prod_{j=1}^qg_j\Big)\|\mc A\|,
    \end{align*}
    for any $q'\in[q]$.
\end{lemma}
A useful corollary can be obtained by replacing $\mc A$ in \lem{commu-bound} with a term from a $k$-local superoperator.
\begin{corollary}
\label{cor:commu-bound-sum}
    Let $\mc K_1, \ldots, \mc K_q$ be a sequence of superoperators on a size-$N$ lattice $\Lambda$ such that $\mc K_j = \sum_{v=1}^{M} \mc K_{j,v}$ is $k$-local and $g_j$-extensive for all $j\in [q  ]$. Then we have 
    \begin{align*}
        \sum_{v_1, \ldots, v_q=1}^{M} \big\|[\mc K_{1,v_1}, \ldots, \mc K_{q,v_q}]\big\|  \le N(q-1)!(2k)^{q-1}\prod_{j=1}^q g_j,
    \end{align*}
\end{corollary}
\begin{proof}
We apply \lem{commu-bound} to the superoperator sequence of length $q-1$, $\mc K_1, \ldots, \mc K_{q-1}$, and set the fixed operator $\mc A = \mc K_{q,v_q}$. Summing over the index $v_q$, we obtain
\begin{align*}
\sum_{v_1, \ldots, v_q=1}^{M} \big\|[\mc K_{1,v_1}, \ldots, \mc K_{q,v_q}]\big\| & \le \sum_{v_q=1}^{M} (q-1)! (2k)^{q-1} \Big(\prod_{j=1}^{q-1} g_j\Big) \|\mc K_{q,v_q}\| \\
&= (q-1)! (2k)^{q-1} \Big(\prod_{j=1}^{q-1} g_j\Big) \sum_{v_q=1}^{M}\|\mc K_{q,v_q}\|\\
&\le (q-1)!\Big(\prod_{j=1}^{q-1} 2k g_j\Big)  Ng_q \\
&=N(q-1)!(2k)^{q-1}\prod_{j=1}^q g_j,
\end{align*}
where the last inequality uses the property that $\mc K_q$ is $g_q$-extensive.
\end{proof}

We can now bound the doubly right-nested commutators defined in Eq.~\eq{def-beta-bar-general} using \lem{commu-bound}. In the following theorem, we turn the locality and extensiveness assumptions into the commutator estimate needed by \thm{general-bch-truncation}.
\begin{theorem}\label{thm:doubly-nested-bound}
    Let $\mathcal{K} = \sum_{v=1}^{M} \mathcal{K}_{v}$ be a $k$-local and $g$-extensive superoperator on a lattice $\Lambda$, and let $\mathcal{A}$ be a superoperator with $|\supp(\mc A)|\le k$. For a given sequence of positive integers $q_1, q_2, \ldots, q_d$, let $P_{r} = 1+\sum_{j=r}^{d} q_j$ for $r=1, 2, \ldots, d$ and $P_{d+1}=1$. Then the following inequality holds for any position index $q' \in [q_d]$
    \begin{align}
    \label{eq:double-comm-k-local}
        \begin{multlined}
            \sum_{v_{1,1}, \ldots, v_{d, q_d}=1}^{M} \bigg\| \Big[ [\mathcal{K}_{v_{1,1}}, \ldots, \mathcal{K}_{v_{1, q_1}}], \ldots, [\mathcal{K}_{v_{d-1,1}}, \ldots, \mathcal{K}_{v_{d-1, q_{d-1}}}], \\
        [\mathcal{K}_{v_{d,1}}, \ldots, \mathcal{K}_{v_{d, q'}}, \mathcal{A}, \mathcal{K}_{v_{d, q'+1}}, \ldots, \mathcal{K}_{v_{d, q_d}}] \Big] \bigg\| 
        \end{multlined}\le \bigg(\prod_{r=1}^{d}P_{r+1} q_r!(2kg)^{q_r}\bigg)\|\mc A\|.
    \end{align}
\end{theorem}

\begin{proof}
    We first decompose each generator $\mathcal{K}_{v}$ into a sum of $k$-local terms. Specifically, let
    \begin{align*}
        \mc K_v = \sum_{\gamma=1}^{\Gamma} \mc K_{\gamma}^{(v)},
    \end{align*}
    where each $\mc K_{\gamma}^{(v)}$ acts non-trivially on at most $k$ sites. To simplify the sum, we re-index these terms as a single sequence $\hat{\mathcal{K}}_{j}$ for $j=1, \ldots, M\Gamma$, such that
    \begin{align*}
        \mc K = \sum_{v=1}^M \mc K_v = \sum_{j=1}^{M\Gamma} \hat{\mc K}_{j}.
    \end{align*}
    By the linearity of the commutator and the triangle inequality, the left-hand side of Eq.~\eqref{eq:double-comm-k-local} is bounded by the sum of the norms of the commutators of the constituent $k$-local terms
    \begin{align}
        \label{eq:triangle-ineq-step}
        \begin{multlined}
             \sum_{v_{1,1}, \ldots, v_{d, q_d}=1}^{M} \bigg\| \Big[ [\mathcal{K}_{v_{1,1}}, \ldots, \mathcal{K}_{v_{1, q_1}}], \ldots, 
        [\mathcal{K}_{v_{d,1}}, \ldots, \mathcal{K}_{v_{d, q'}}, \mathcal{A}_X, \ldots, \mathcal{K}_{v_{d, q_d}}] \Big] \bigg\| \\
        \le \sum_{u_{1,1}, \ldots, u_{d, q_d}=1}^{M\Gamma} \bigg\| \Big[ [\hat{\mathcal{K}}_{u_{1,1}}, \ldots, \hat{\mathcal{K}}_{u_{1, q_1}}], \ldots, 
        [\hat{\mathcal{K}}_{u_{d,1}}, \ldots, \hat{\mathcal{K}}_{u_{d, q'}}, \mathcal{A}, \ldots, \hat{\mathcal{K}}_{u_{d, q_d}}] \Big] \bigg\|.
        \end{multlined}
    \end{align}
    For any layer $r \in [d]$, let $\vec{u}_r = (u_{r,1}, \ldots, u_{r, q_r})$ denote the vector of indices ranging from $1$ to $M\Gamma$. We define the nested commutator for layer $r$ as
    \begin{align*}
        \hat{\mathcal{K}}_{\vec{u}_r} := [\hat{\mathcal{K}}_{u_{r,1}}, \ldots, \hat{\mathcal{K}}_{u_{r, q_r}}].
    \end{align*}
    Similarly, let $\hat{\mathcal{W}}_{\vec{u}_d} := [\hat{\mathcal{K}}_{u_{d,1}}, \ldots, \mathcal{A}_X, \ldots, \hat{\mathcal{K}}_{u_{d, q_d}}]$ denote the innermost term containing $\mc A$ at the $q'$-th position.
    The RHS of Eq.~\eqref{eq:triangle-ineq-step} can be rewritten as
    \begin{align*}
        \sum_{\vec{u}_{1}, \ldots, \vec{u}_{d}} \bigg\| \Big[ \hat{\mathcal{K}}_{\vec{u}_1}, \ldots, \hat{\mathcal{K}}_{\vec{u}_{d-1}}, 
        \hat{\mathcal{W}}_{\vec{u}_d} \Big] \bigg\|.
    \end{align*}
    Define the partial nested commutator starting from layer $r$ as
    \begin{align*}
        \mathcal{W}_{\vec{u}_{r}, \ldots, \vec{u}_d} := \Big[ \hat{\mathcal{K}}_{\vec{u}_r}, \ldots, \hat{\mathcal{K}}_{\vec{u}_{d-1}}, \hat{\mathcal{W}}_{\vec{u}_d} \Big].
    \end{align*}
    We prove the following bound by backward induction on $r$ from $d$ to $1$
    \begin{align} \label{eq:induct-W}
        \sum_{\vec{u}_{r}, \ldots, \vec{u}_{d}} \Big\|\mathcal{W}_{\vec{u}_{r}, \ldots, \vec{u}_d}\Big\|\le \bigg(\prod_{j=r}^{d}P_{j+1} q_j!(2kg)^{q_j}\bigg)\|\mc A\|.
    \end{align}
    
    Consider the base case $r=d$. The term is $\sum_{\vec{u}_d} \|\hat{\mathcal{W}}_{\vec{u}_d}\|$. Since $\mathcal{W}_{\vec{u}_d}$ involves one fixed operator $\mathcal{A}$ and $q_d$ operators from a $k$-local set, we apply \lem{commu-bound}. This yields the bound $q_d!(2kg)^{q_d}\|\mathcal{A}\|$. Since $P_{d+1}=1$, the base case holds.

    Assume Eq.~\eqref{eq:induct-W} holds for $r$. We consider the sum for $r-1$
    \begin{align*}
        \sum_{\vec{u}_{r-1}, \ldots, \vec{u}_d} \Big\|\mathcal{W}_{\vec{u}_{r-1}, \ldots, \vec{u}_d}\Big\| 
        &= \sum_{\vec{u}_{r}, \ldots, \vec{u}_d} \sum_{\vec{u}_{r-1}} \bigg\| \Big[ \hat{\mathcal{K}}_{\vec{u}_{r-1}}, \mathcal{W}_{\vec{u}_{r}, \ldots, \vec{u}_d} \Big] \bigg\|.
    \end{align*}
    Let $X_{\mathcal{W}} = \supp(\mathcal{W}_{\vec{u}_{r}, \ldots, \vec{u}_d})$. Its size is bounded by $|X_{\mathcal{W}}| \le k P_r$. The inner commutator is nonzero only if at least one operator in $\hat{\mathcal{K}}_{\vec{u}_{r-1}}$ overlaps with $X_{\mathcal{W}}$.
    We explicitly expand the sum over $\vec{u}_{r-1}$ by iterating over which operator overlaps with $X_{\mathcal{W}}$. Let $j \in [q_{r-1}]$ be the index of this operator. We can upper bound the sum by
    \begin{align*}
        & \sum_{\vec{u}_{r-1}} \bigg\| \Big[ \hat{\mathcal{K}}_{\vec{u}_{r-1}}, \mathcal{W}_{\vec{u}_{r}, \ldots, \vec{u}_d} \Big] \bigg\| \\
        \le~& \sum_{j=1}^{q_{r-1}} \sum_{\substack{u_{r-1, j} \in [M\Gamma] \\ \supp(\hat{\mathcal{K}}_{u_{r-1, j}}) \cap X_{\mathcal{W}} \neq \emptyset}} \sum_{\substack{u_{r-1, k} \in [M\Gamma] \\ (\forall k \neq j)}} \Big\| \big[ \hat{\mathcal{K}}_{u_{r-1,1}}, \ldots, \hat{\mathcal{K}}_{u_{r-1, q_{r-1}}}, \mathcal{W}_{\vec{u}_{r}, \ldots, \vec{u}_d} \big] \Big\| \\
        \le~& \sum_{j=1}^{q_{r-1}} \sum_{\substack{u_{r-1, j} \in [M\Gamma] \\ \supp(\hat{\mathcal{K}}_{u_{r-1, j}}) \cap X_{\mathcal{W}} \neq \emptyset}} (q_{r-1}-1)! (2kg)^{q_{r-1}-1} \|\hat{\mathcal{K}}_{u_{r-1, j}}\| \cdot 2 \Big\|\mathcal{W}_{\vec{u}_{r}, \ldots, \vec{u}_d}\Big\| \\
        \le~& \sum_{j=1}^{q_{r-1}} (k P_r g) \cdot (q_{r-1}-1)! (2kg)^{q_{r-1}-1} \cdot 2 \Big\|\mathcal{W}_{\vec{u}_{r}, \ldots, \vec{u}_d}\Big\| \\
        =~& q_{r-1} \cdot (k P_r g) \cdot (q_{r-1}-1)! (2kg)^{q_{r-1}-1} \cdot 2 \Big\|\mathcal{W}_{\vec{u}_{r}, \ldots, \vec{u}_d}\Big\| \\
        =~& P_r \cdot q_{r-1}! \cdot (2kg)^{q_{r-1}} \Big\|\mathcal{W}_{\vec{u}_{r}, \ldots, \vec{u}_d}\Big\|.
    \end{align*}
    Here, the second inequality follows from applying \lem{commu-bound} to the sum over the $q_{r-1}-1$ operators not at index $j$ with $\mc A:= \hat{\mathcal{K}}_{u_{r-1, j}}$.  The third inequality uses the $g$-extensiveness property: summing the norm of $\hat{\mathcal{K}}_{u_{r-1, j}}$ over all operators overlapping with $X_{\mathcal{W}}$ yields at most $|X_{\mathcal{W}}|g \le k P_r g$.
    
    Substituting this back and using the inductive hypothesis leads to
    \begin{align*}
        \sum_{\vec{u}_{r-1}, \ldots, \vec{u}_d} \Big\|\mathcal{W}_{\vec{u}_{r-1}, \ldots, \vec{u}_d}\Big\| 
        &\le P_r q_{r-1}! (2kg)^{q_{r-1}} \sum_{\vec{u}_{r}, \ldots, \vec{u}_d} \Big\|\mathcal{W}_{\vec{u}_{r}, \ldots, \vec{u}_d}\Big\| \\
        &\le P_r q_{r-1}! (2kg)^{q_{r-1}}  \bigg(\prod_{j=r}^{d}P_{j+1} q_j!(2kg)^{q_j}\bigg)\|\mc A\|.
    \end{align*}
    This completes the induction.
\end{proof}

Plugging this bound into \thm{general-bch-truncation}, we can bound the truncation error for a $k$-local system as follows. Note that if all $\mathcal{K}_v$ are Hamiltonians, this bound matches the one in \citet[Theorem 3]{mizuta2025commutator}. However, we use a distinct proof that applies to general superoperators beyond lattice systems. The following lemma records the resulting truncation guarantee in a form that will be used in the extrapolation analysis.
\begin{lemma}
\label{lem:bch-approx}
    Given $k$-local and $g$-extensive superoperators $\mc K_v$, for any $\eps\in (0,1/2)$ and  $q_0 \ge \log(2N/\eps)$, if $4e^2kgq_0\le 1$, then we have \begin{align*}
        \left\|\exp\Big(\sum_{q=1}^{q_0} \Phi_q\Big) - \prod_{v = 1}^Me^{\mc K_v} \right\|\le \eps.
    \end{align*}
\end{lemma}
\begin{proof} 
    For any positive integer sequence $q_1, \ldots, q_{d+1}$, let 
    \begin{align*}
        P_{i} = \sum_{j=i}^{d+1} q_j,
    \end{align*}
    for $i\le d+1$, $P_{d+2} = 1$, and $q = q_1+\cdots+q_{d+1}$. Then, we have 
    \begin{align*}
    &\bar{\alpha}_{\comm}^{(q_1, \ldots, q_{d+1})}\\
    =~& \sum_{v_{d+1, q_{d+1}}=1}^{M} \sum_{v_{1,1}, \ldots, v_{d+1, q_{d+1}-1} = 1}^M \bigg\| \Big[ 
         [\mc K_{v_{1,1}}, \ldots, \mc K_{v_{1,q_1}}], \ldots, 
        [\mc K_{v_{d+1,1}}, \ldots, \mc K_{v_{d+1,q_{d+1}}}]
     \Big] \bigg\| \\
    \le~& \bigg(\prod_{i=1}^{d+1}P_{i+1}q_i!(2kg)^{q_i }\bigg)\frac{1}{g}\sum_{v_{d+1, q_{d+1}}=1}^{M} \|\mc K_{v_{d+1, q_{d+1}}}\|\\
    =~&  \bigg(\prod_{i=1}^{d+1}P_{i+1}q_i!(2kg)^{q_i }\bigg)N,
\end{align*}
where the third line follows from \thm{doubly-nested-bound}, and the fourth line follows from 
\begin{align*}
    \sum_{v_{d+1,q_{d+1}}=1}^{M}\|\mc K_{v_{d+1, q_{d+1}}}\|\le \sum_{x\in\Lambda}\sum_{\substack{v = 1, \ldots, M \\ \supp(\mc K_v)\ni x}} \|\mc K_{v}\|\le gN.
\end{align*}
For $q_1, \ldots, q_{d+1}\le q_0$, we have 
\begin{align*}
    \bar{\alpha}_{\comm}^{(q_1, \ldots, q_{d+1})}\le q^{d}\bigg(\prod_{i=1}^{d+1}q_0^{q_{i}}(2kg)^{q_i }\bigg)N = q^{d} (2q_0kg)^q N,
\end{align*}
where the inequality follows from $P_{i+1}\le q$ (noting that $P_{d+2} = 1$ does not contribute a factor of $q$). 
Then, we have 
\begin{align*}
    \bar{\alpha}_{\comm,q_0} &\le (2q_0 kg)^{q_0+1} N+ \sum_{r = 0}^{\infty}\frac{q^{r} }{(r+1)!} \sum_{q=q_0+2}^{\infty}\sum_{\substack{1\le p_1, \ldots, p_{r+1}\le q_0\\p_1+\cdots+p_{r+1} = q}}(2q_0 kg)^q N \\ 
    &\le  (2q_0 kg)^{q_0+1} N+  \sum_{q=q_0+2}^{\infty}\sum_{\substack{1\le p_1, \ldots, p_{r+1}\le q_0\\p_1+\cdots+p_{r+1} = q}} (2eq_0 kg)^q N\\
    &\le(2q_0 kg)^{q_0+1} N+  \sum_{q=q_0+2}^{\infty} (4eq_0 kg)^q N   \\
    &\le (2q_0kg)^{q_0+1} N + (4eq_0 kg)^{q_0+1} N \sum_{q = 1}^{\infty} 2^{-q} \\
    &\le 2(4eq_0kg)^{q_0+1}N \le (4e^2q_0 kg)^{q_0+1}\le 1,
\end{align*}
where the third line follows from 
\begin{align*}
    \sum_{\substack{1\le q_1, \ldots, q_{r+1}\le q_0\\q_1+\cdots q_{r+1}=q}}1\le  \sum_{\substack{1\le q_1, \ldots, q_{r+1}\\q_1+\cdots q_{r+1}=q}}1=\binom{q-1}{r}\le 2^q,
\end{align*}
and the fourth line follows from $e^{-q_0-1} \le \eps/(2N) \le 1/(2N)$.
As $\bar{\alpha}_{\comm,q_0}\le 1$ is satisfied, we can apply \thm{general-bch-truncation} to obtain 
\begin{align*}
\left\|\exp\Big(\sum_{q=1}^{q_0} \Phi_q\Big) - \prod_{v = 1}^Me^{\mc K_v} \right\|
\le  e\bar{\alpha}_{\comm,q_0}
\le 2e(4eq_0kg)^{q_0+1}N\le 2e^{-q_0}N\le \eps.
\end{align*}
\end{proof}
Based on the convergence condition in Eq.~\eq{converg-bch} and the bound of $\alpha_{\rm comm}$ in Eq.~\eq{trivial-alpha}, the BCH series $\sum_{q=1}^{\infty} \Phi_q$ is guaranteed to converge only when $\sum_{v=1}^M \|\mc K_v\| = O(1)$, which requires $g = O(1/N)$. In contrast, \lem{bch-approx}  shows that the exponential of the truncated BCH formula approximates the exponential product in a much larger regime of $g = O(1/(k\log N))$, where the BCH series is not guaranteed to converge. 

%%%%%%%%%%%%%%%%%%%%%%%%%%%%%%%%%%%%%%%%%%%%%%%%%%%%%%%%%%%%%%%%%%%%%%%%%%%%%%%%%%%%%%%%%%%%%%%%%%%%%%%%%%%%%%%%

\subsubsection{Application to second-quantized electronic structure}
Following the notation of \citet{childs2021theory}, we consider the second-quantized electronic structure Hamiltonian in the plane-wave dual basis. In this representation, the Hamiltonian for a system with $n$ spin-orbitals takes the form
\begin{align}
    \label{eq:def-tuv}
    H = \underbrace{\sum_{j,k=1}^n t_{j,k} A_j^\dagger A_k}_{T} 
      + \underbrace{\sum_{j=1}^n u_j N_j}_{U} 
      + \underbrace{\sum_{j,k=1}^n g_{j,k} N_j N_k}_{V},
\end{align}
where $A_j^\dagger, A_j$ are the creation and annihilation operators satisfying $\{A_j, A_k^\dagger\} = \delta_{jk}$, and $N_j = A_j^\dagger A_j$. Here, $T$ represents the kinetic energy with hopping amplitudes $t_{j,k}$, $U$ denotes the single-particle potential with coefficients $u_j$, and $V$ describes the two-body interaction with strength $g_{j,k}$.
Note that $\|A_j^{\dagger}\|=\|A_k\|=\|N_j\|=1$. We can apply the triangle inequality to upper bound the norms $\|T\|$, $\|U\|$, and $\|V\|$ by the $1$-norm of the corresponding coefficient vectors $\|\vec{t}\|_1$, $\|\vec{u}\|_1$, and $\|\vec{v}\|_1$. \citet{childs2021theory} proved the follow upper bounds for the $1$-norm and $\infty$-norm of the coefficient vectors. 
\begin{lemma}[{\cite[Proposition F.2]{childs2021theory}}]
\label{lem:norm-tuv}
    The $1$-norm and $\infty$-norm of the coefficients vectors of the kinetic operator $T$ and potential operators $U$, $V$ in Eq.~\eq{def-tuv} can be bounded as \begin{align*}
        &\|\vec{t}\|_{\infty} = O(1/n), && \|\vec{t}\|_1 = O(n) \\
        &\|\vec{u}\|_{\infty} = O(n), && \|\vec{u}\|_1 = O(n^2) \\
        &\|\vec{g}\|_{\infty} = O(1), && \|\vec{g}\|_1  = O(n^2). 
    \end{align*}
\end{lemma}
Denote the summand in Eq.~\eq{def-tuv} by $H_{v}$, which gives \begin{align*}
    H = \sum_{v = 1}^{M} H_{v},
\end{align*}
with $M = 2n^2 +n$.
The right-nested commutators of $H_v$ can be expressed as \begin{align*}
    W = \sum_{\vec{j}, \vec{k},\vec{\ell}}w_{\vec{j}, \vec{k},\vec{\ell}}\cdots (A_{j_x}^{\dagger} A_{k_x}) \cdots (N_{\ell_y})\cdots,
\end{align*}
where $\vec{j}, \vec{k},\vec{\ell}$ are vectors of orbitals. The largest number of $A_{j_x}^{\dagger} A_{k_x}$ and $N_{\ell_y}$ terms in the summand is called the ``layer'' of $W$. For any $W$, the norms of the commutators of terms in $T$, $U$, and $V$ with $W$ are bounded by the $\infty$-norms of $t$, $u$, and $g$, and the $1$-norm of $w$.
\begin{lemma}[{\cite[Proposition F.4]{childs2021theory}}]
\label{lem:recursion-W}
     Let an electronic-structure Hamiltonian be given as in Eq.~\eq{def-tuv}. The following statements hold for a general second-quantized operator $W$ with $q$ layers:
     \begin{enumerate}
         \item $\tilde{W}=[T, W]$ is an operator with $q$ layers and $\|\vec{\tilde{w}}\|_1 \leq 2 q n\|\vec{t}\|_{\infty}\|\vec{w}\|_1 $;
         \item $\tilde{W}=[U, W]$ is an operator with $q$ layers and $\|\vec{\tilde{w}}\|_1 \leq 2 q\|\vec{u}\|_{\infty}\|\vec{w}\|_1$; and 
         \item $\tilde{W}=[V, W]$ is an operator with $q+1$ layers and $\|\vec{\tilde{w}}\|_1 \leq 4 q n\|\vec{g}\|_{\infty}\|\vec{w}\|_1$.
     \end{enumerate}
\end{lemma}
A Corollary of this lemma is an upper bound for the sum of right-nested commutator norms. 
\begin{corollary}
\label{cor:right-comm-second}
    The sum of the norms of the right-nested commutators of $H_v$ satisfies 
    \begin{align*}
        \sum_{v_1, \ldots, v_q=1}^{M} \|[H_{v_1}, \ldots, H_{v_q}]\| = \mathcal{O}(q! n^q),
    \end{align*}
    for any $q\ge 2$. In addition, each $[H_{v_1}, \ldots, H_{v_q}]$ has at most $q$ layers. 
\end{corollary}
\begin{proof}
    We prove the statement by induction on $q$. 
    
    Consider the base case $q=2$. The LHS is $\sum_{v_1, v_2=1}^{M} \|[H_{v_1}, H_{v_2}]\|$, where each $H_{v_1}$ and $H_{v_2}$ is a term belonging to $T$, $U$, or $V$. Note that terms from $U$ and $V$ commute with each other, so we only need to consider commutators between $T$ and $\{U,V\}$. This yields
    \begin{align*}
        \sum_{v_1, v_2=1}^{M} \|[H_{v_1}, H_{v_2}]\| 
        &= \mathcal{O}\big(n\|\vec{t}\|_{\infty} \|\vec{u}\|_1+ n\|\vec{t}\|_{\infty} \|\vec{g}\|_1 + n\|\vec{t}\|_{\infty} \|\vec{t}\|_1\big) \\
        &= \mathcal{O}(n^2).
    \end{align*}
    The first equality follows from \lem{recursion-W} and the second follows from the coefficient norms in \lem{norm-tuv}. In addition, the layer of $[H_{v_1}, H_{v_2}]$ is at most $2$. Therefore, the base case holds. 
    
    Assume the statement holds for $q-1$. Consider the sum for $q$: 
    \begin{align*}
         \sum_{v_1, \ldots, v_q=1}^{M} \|[H_{v_1}, \ldots, H_{v_q}]\| = \sum_{v_1=1}^M  \sum_{v_2, \ldots, v_q=1}^{M} \Big\| \big[H_{v_1}, [H_{v_2}, \ldots, H_{v_q}] \big] \Big\|. 
    \end{align*} 
    By the induction hypothesis, $[H_{v_2}, \ldots, H_{v_q}]$ has at most $q-1$ layers.
    According to \lem{recursion-W}, commuting with terms belonging to $T$, $U$, or $V$ increases the norm by a factor proportional to $(q-1)n\|\vec{t}\|_\infty = \mathcal{O}(q)$, $(q-1)\|\vec{u}\|_\infty =\mathcal{O}(qn)$, or $(q-1)n\|\vec{g}\|_\infty=\mathcal{O}(qn)$, respectively. Summing over all possible $H_{v_1}$, we have
    \begin{align*}
        \sum_{v_1=1}^M \sum_{v_2, \ldots, v_q=1}^{M} \Big\| \big[H_{v_1}, [H_{v_2}, \ldots, H_{v_q}] \big] \Big\| 
        = \mathcal{O}\Big(qn \sum_{v_2, \ldots, v_q=1}^{M} \|[H_{v_2}, \ldots, H_{v_q}]\|\Big) 
        = \mathcal{O}(q!n^{q}),
    \end{align*}
    where the last equality follows from the induction hypothesis. By \lem{recursion-W}, commuting with an $H_v$ increases the layer by at most one, and hence $[H_{v_1}, \ldots, H_{v_q}]$ has at most $q$ layers.
\end{proof}



Let $\mc{K}_v:=-\i \ad_{H_v}$ and \begin{align}
\label{eq:def-K-second-quantized}
    \mc{K}:= \sum_{v=1}^M  \mc K_v= -\i \ad_H. 
\end{align}
We now apply \thm{general-bch-truncation} to bound the truncation error of the BCH formula of $\prod_{v=1}^M\mc K_v$. The following lemma gives the required doubly right-nested commutator bound for the second-quantized electronic-structure Hamiltonian.
\begin{lemma}
    Let $\mc K = \sum_{v=1}^M \mc K_v$ be the Liouvillian in Eq.~\eq{def-K-second-quantized}. For a given sequence of positive integers $q_1, q_2, \ldots, q_d$, let $P_{r} = \sum_{j=r}^{d} q_j$ for $r=1, 2, \ldots, d$ and $P_{d+1}=1$. Then, the doubly right-nested commutator bound $\bar{\alpha}_{\comm}^{(q_1, \ldots, q_{d})}$ for $\mc K$ (see Eq.~\eq{def-beta-bar-general}) satisfies
    \begin{align}
    \label{eq:double-comm-second-quantized}
        \bar{\alpha}_{\comm}^{(q_1, \ldots, q_{d})}= O\bigg(\prod_{r=1}^{d}P_{r+1} q_r!n^{q_r}\bigg).
    \end{align}
\end{lemma}
\begin{proof}
    First, we reduce the problem to bounding commutators of Hamiltonian terms. By applying the Jacobi identity $[\ad_{X}, \ad_Y] = \ad_{[X,Y]}$ recursively and using $\|\ad_X\| \le 2\|X\|$, the doubly right-nested commutator bound of $\mc K_v$ satisfies
    \begin{align*}
    \bar{\alpha}_{\comm}^{(q_1, \ldots, q_{d})}
    &= \sum_{v_{1,1}, \ldots, v_{d, q_{d}} = 1}^M \bigg\| \ad_{\Big[ 
         [H_{v_{1,1}}, \ldots, H_{v_{1,q_1}}], \ldots, 
        [H_{v_{d,1}}, \ldots, H_{v_{d,q_{d}}}]
     \Big]} \bigg\|\\
     &\le 2\sum_{v_{1,1}, \ldots, v_{d, q_{d}} = 1}^M \bigg\|{\Big[ 
         [H_{v_{1,1}}, \ldots, H_{v_{1,q_1}}], \ldots, 
        [H_{v_{d,1}}, \ldots, H_{v_{d,q_{d}}}]
     \Big]} \bigg\|.
    \end{align*}
    Therefore, it suffices to bound the sum of doubly right-nested commutators of $H_v$. 
    
    To simplify the notation, for any layer $r \in [d]$, let $\vec{v}_r = (v_{r,1}, \ldots, v_{r, q_r})$ denote the index vector, and define the nested commutator for that layer as
    \begin{align*}
        H_{\vec{v}_r} := [H_{v_{r,1}}, \ldots, H_{v_{r, q_r}}].
    \end{align*}
    We define the partial nested commutator starting from layer $r$ as
    \begin{align*}
        W_{\vec{v}_{r}, \ldots, \vec{v}_d} := \Big[ H_{\vec{v}_r}, \ldots, H_{\vec{v}_{d-1}}, H_{\vec{v}_d} \Big].
    \end{align*}
    We prove the following bound by backward induction on $r$ from $d$ to $1$:
    \begin{align} \label{eq:induct-W-elec}
        \sum_{\vec{v}_{r}, \ldots, \vec{v}_{d}} \Big\|W_{\vec{v}_{r}, \ldots, \vec{v}_d}\Big\| = \mathcal{O}\bigg(\prod_{j=r}^{d}P_{j+1} q_j! n^{q_j}\bigg).
    \end{align} In addition, we show that $W_{\vec{v}_{r}, \ldots, \vec{v}_d}$ has at most $P_r$ layers.  
    
    Consider the base case $r=d$. By \cor{right-comm-second}, the LHS of Eq.~\eq{induct-W-elec} is \begin{align*}
        \sum_{\vec{v}_d} \|H_{\vec{v}_d}\| = \mc{O}(q_d ! n^{q_d}) = \mc{O}(P_{d+1}q_d ! n^{q_d}),
    \end{align*} and $H_{\vec{v}_d}$ has $q_d$ layers, so the base case holds. 

    Assume Eq.~\eqref{eq:induct-W-elec} holds for $r$. We consider the sum for $r-1$:
    \begin{align*}
        \sum_{\vec{v}_{r-1}, \ldots, \vec{v}_d} \Big\|W_{\vec{v}_{r-1}, \ldots, \vec{v}_d}\Big\| 
        &= \sum_{\vec{v}_{r}, \ldots, \vec{v}_d} \sum_{\vec{v}_{r-1}} \bigg\| \Big[ H_{\vec{v}_{r-1}}, W_{\vec{v}_{r}, \ldots, \vec{v}_d} \Big] \bigg\|.
    \end{align*}
    If $q_{r-1}\ge 2$, we have \begin{align*}
        \sum_{\vec{v}_{r}, \ldots, \vec{v}_d} \sum_{\vec{v}_{r-1}} \bigg\| \Big[ H_{\vec{v}_{r-1}}, W_{\vec{v}_{r}, \ldots, \vec{v}_d} \Big] \bigg\| &\le 2\Big(\sum_{\vec{v}_{r-1}}\|H_{\vec{v}_{r-1}}\|\Big)\Big(\sum_{\vec{v}_{r}, \ldots, \vec{v}_d}\|W_{\vec{v}_{r}, \ldots, \vec{v}_d}\|\Big) \\
        &= \mc{O}\Big(q_{r-1}!n^{q_{r-1}}\prod_{j=r}^{d}P_{j+1} q_j! n^{q_j}\Big)\\
        &= \mc{O}\Big(\prod_{j=r-1}^{d}P_{j+1} q_j! n^{q_j}\Big),
    \end{align*}
    where the second line follows from \cor{right-comm-second} and the induction hypothesis. The layer count of $W_{\vec{v}_{r-1}, \ldots, \vec{v}_d}$ is the sum of the layer counts of $H_{\vec{v}_{r-1}}$ and $W_{\vec{v}_{r}, \ldots, \vec{v}_d}$, which is $q_{r-1}+ P_{r} = P_{r-1}$. 
    
    Then we consider the case of $q_{r-1}=1$. We have 
    \begin{align*}
        \sum_{\vec{v}_{r-1}, \ldots, \vec{v}_d} \Big\|W_{\vec{v}_{r-1}, \ldots, \vec{v}_d}\Big\| 
       = \sum_{v}\sum_{\vec{v}_{r}, \ldots, \vec{v}_d} \Big\|[H_{v},W_{\vec{v}_{r}, \ldots, \vec{v}_d}]\Big\|.
    \end{align*}
    By the induction hypothesis, $W_{\vec{v}_{r}, \ldots, \vec{v}_d}$ has at most $P_{r}$ layers. Using the same argument in the proof of \cor{right-comm-second}, commuting with an $H_v$ increases the norm by a factor of $\mc{O}(P_r n)$. Therefore, we have \begin{align*}
        \sum_{v}\sum_{\vec{v}_{r}, \ldots, \vec{v}_d} \Big\|[H_{v},W_{\vec{v}_{r}, \ldots, \vec{v}_d}]\Big\| &= \mc{O}\Big(P_r n \sum_{\vec{v}_{r}, \ldots, \vec{v}_d} \|W_{\vec{v}_{r}, \ldots, \vec{v}_d}\|\Big)\\
        &= \mc{O}\Big(P_r n \prod_{j=r}^{d}P_{j+1} q_j! n^{q_j}\Big) \\
        &= \mc{O}\Big( \prod_{j=r-1}^{d}P_{j+1} q_j! n^{q_j}\Big)
    \end{align*}
    where the second line follows from the induction hypothesis. By \lem{recursion-W}, $W_{\vec{v}_{r-1}, \ldots, \vec{v}_d}$ has at most $1+P_{r} = P_{r-1}$ layers.  
    This completes the induction. Setting $r=1$ yields the lemma statement.
\end{proof}


\section{Analysis of the extrapolation error}
\textbf{Roadmap.} In this section, we derive the 1D joint error mitigation scheme and analyze the extrapolation error. Specifically, we upper bound the extrapolation error by the linear combination of the simulation errors of each noisy quantum circuit, which can be further decomposed into the BCH truncation error of the noisy quantum channel plus the distance between the truncated BCH formula and the ideal evolution channel.

Given a Hamiltonian $H = \sum_{\gamma=1}^{\Gamma} H_{\gamma}$, an initial state $\rho_0$, and an observable $O$, we aim to estimate \begin{align*}
    \operatorname{Tr}[Oe^{-\i \ad_H T} \rho_0],
\end{align*}
for any simulation time $T$.
Suppose that the Hamiltonian evolution driven by each $H_{\gamma}$ can be efficiently implemented by a layer of gates. Then, for sufficiently small $t$, we can implement the Hamiltonian evolution $e^{-\i H t}$ approximately by a $p$-th order product formula of the form
\begin{align*}
    P(t) = \prod_{v=1}^Me^{-\i t a_{v} H_{\gamma_v}},
\end{align*}
where the product is ordered from right to left. Here, $a_{v}$ are real coefficients, $M = \Upsilon\Gamma$ is the total number of unitary gates, and the sequence $\gamma_v$ specifies which Hamiltonian is applied at step $v$. Denote $a_{\max} = \max_{v\in [M]} |a_v|$. The corresponding ideal channel $\P(t)$ is a superoperator given by
\begin{align*}
    \P(t) = \prod_{v=1}^Me^{-\i t a_{v} \ad_{H_{\gamma_v}}},
\end{align*}
where $\A_{\gamma}:=\ad_{H_{\gamma}}$ denotes the adjoint action of $H_{\gamma}$. 

Suppose that the system is subject to physical noise with a controllable strength $\lambda$. Specifically, after each layer of unitary gates, a noise channel $e^{\lambda \mc L_{v}}$ is applied, where $\mc L_{v}$ is a Lindbladian. The resulting noisy channel is
\begin{align*}
    \tilde{\P}(t,\lambda)=\prod_{v=1}^M\left( e^{\lambda \mc L_{v}} e^{-\i t a_{v} \ad_{H_{\gamma_v}}} \right).
\end{align*} 
% In experimental implementations, the noise strength is lower-bounded by the device's intrinsic noise level, denoted by $\lambda_{\min}$. \xz{add some overview on how to tune the noise level like unitary folding or learning and recovering the Pauli channel.} \sz{Folding not feasible in near-term, since the circuit depth beyond allowance.} Therefore, we assume that $\lambda$ can be varied within the range $[\lambda_{\min}, \infty)$. \sz{In practice, too large $\lambda$ may cause the experiment 'collapse'.}

In experimental implementations, the noise strength is lower-bounded by the device's intrinsic noise level, denoted by $\lambda_{\min}$. To access larger noise strengths, one common approach is \textit{unitary folding}~\cite{giurgica-tiron_digital_2020}, which replaces each unitary gate $U$ by $U(U^\dagger U)^n$, effectively scaling the noise by a factor of $1+2n$ without requiring explicit knowledge of the noise model. However, folding increases the circuit depth proportionally, which may exceed the coherence budget of near-term devices.
Alternatively, if the noise channel can be characterized---for instance, learned as a sparse Pauli--Lindblad model~\cite{berg_probabilistic_2023}---one can probabilistically insert recovery gates to directly emulate amplified noise at the pulse level without additional circuit depth~\cite{li2017efficient}. In practice, the noise strength cannot be amplified arbitrarily without degrading the signal, but moderate amplification is experimentally feasible. We therefore assume that $\lambda$ can be tuned within the range $[\lambda_{\min}, \infty)$ for the purpose of asymptotic analysis.
Given $s\in (0,1)$ such that $1/s\in \mathbb{Z}$, we implement $\mc P(sT)^{1/s}$ with initial state $\rho_0$ and then measure the observable $O$ to approximate $\operatorname{Tr}[Oe^{-\i \ad_H T}\rho_0]$. We denote the expectation of the measurement results with noisy circuit by  \begin{align*}
    g(s,\lambda) := \operatorname{Tr}[O\tilde{\mc P}(sT,\lambda)^{1/s}\rho_0],
\end{align*}
and note that \begin{align*}
    \lim_{s,\lambda\to 0} g(s,\lambda) = \operatorname{Tr}[Oe^{-\i \ad_H T}\rho_0]
\end{align*}
is the quantity we want. To obtain this limit value, one could perform two successive extrapolations: first, use ZNE to extrapolate $\lambda \to 0$ for a fixed $s$, and then use Trotter error mitigation \cite{watson2025exponentially} to extrapolate $s \to 0$ using the ZNE output. Although this method is possible, it involves extrapolating two independent parameters, which requires more sample points. Furthermore, each extrapolation step amplifies other sources of noise, such as state preparation and measurement (SPAM) errors, making the method less robust.

Therefore, we consider a more efficient and robust approach: a unified 1D extrapolation. By coupling the noise strength and the step size, we reduce this 2D problem to a single-variable extrapolation. In general, we could set $\lambda$ to be an analytic function of $s$, $\lambda(s)$, such that $\lambda(0)=0$. This ensures that the extrapolated function $g(s,\lambda(s))$ has the correct limit value at $s = 0$. Furthermore, its analyticity near $s=0$ guarantees the existence of a power series expansion, which is the fundamental requirement for Richardson extrapolation. For general analytic function $\lambda(s) =\sum_{j=p_0+1}^{\infty} c_j (sT)^j$, the noise error will be dominated by its leading-order term, $O(s^{p_0+1})$. Thus, for analyzing the asymptotic scaling, it is sufficient to consider a simple monomial form  $\lambda(s) = c(sT)^{p_0+1}$ for some $c>0$ and $p_0\in \mb{Z}_+$.

We now determine the optimal exponent $p_0$. The total error of the estimator $g(s, \lambda(s))$ arises from two dominant sources: the algorithmic (Trotter) error, which scales as $O(s^{p+1})$ for a $p$-th order product formula, and the physical noise error, which scales linearly with $\lambda$ in the leading order, and thus scales as $O(s^{p_0+1})$.
The leading-order error of the combined system is governed by the term with the slower decay, i.e., $O(s^{\min\{p+1, p_0+1\}})$. To maximize the efficiency of the extrapolation, we aim to maximize this leading-order exponent. If we choose $p_0 < p$, the noise error dominates ($O(s^{p_0+1})$) and becomes the bottleneck, wasting the precision of the $p$-th order Trotter formula. Conversely, if $p_0 > p$, the Trotter error dominates ($O(s^{p+1})$), implying that we are suppressing the noise strength more aggressively than necessary.
The optimal strategy is to match the scaling of the two error sources by setting $p_0 = p$. This choice ensures that both the leading-order Trotter error and the leading-order noise error appear at the same order, $O(s^{p+1})$, allowing the 1D Richardson extrapolation to eliminate both simultaneously and efficiently. We further remark that this matching strategy is not merely a heuristic. As we will demonstrate in our applications (see \sec{local} and \sec{geo-local}) \xz{add ref to application}, this choice enables our protocol to achieve a simulation depth that saturates the fundamental limits of quantum error mitigation under Pauli noise, as established in Refs.~\cite{wang2021noise, cai2021multi}.

To analyze the Richardson extrapolation error, our goal is to bound the remainder of the power series expansion of \begin{align*}
    f(s):= g(s, c(sT)^{p+1}) =  \operatorname{Tr}[O\tilde{\mc P}(sT,c(sT)^{p+1})^{1/s}\rho_0].
\end{align*}
The $1/s$ exponent in the definition of $f(s)$ makes a direct power series analysis difficult. The standard approach is to first analyze the single-step propagator, $\tilde{\mc P}(sT,c(sT)^{p+1})$, by finding its effective generator. Note that the full channel can be written as \begin{align*}
    \tilde{\mc P}(t,\lambda) = \prod_{v=1}^{2M} e^{\mc K_{v}},
\end{align*} where each $\mc K_v$ is a Lindbladian defined as 
\begin{align}
 	\mathcal{K}_j= \begin{cases}-\i ta_{v}\ad_{H_{\gamma_v}} & \text{if } j = 2v-1\\ \lambda\mathcal{L}_v & \text{if } j = 2v \end{cases}, 	\label{eq:def-K}
\end{align}
for $v \in [M]$. Therefore, the effective generator of $\tilde{\mc P}(t,\lambda)$ is given by the BCH formula
\begin{align*}
    \tilde{\mc P}(t,\lambda)  = \exp\left(\sum_{q=1}^\infty \Phi_q(t,\lambda)\right),
\end{align*}
where $\Phi_q$ is the $q$-th order term in the BCH formula in Eq.~\eq{bch} with $2M$ generators $\mc K_1, \ldots, \mc K_{2M}$ defined in Eq.~\eq{def-K}.
The full evolution is then \begin{align*}
    \tilde{\mc P}(sT,c(sT)^{p+1})^{1/s} = \exp(\frac{1}{sT}\sum_{q=1}^\infty \Phi_q(sT,c(sT)^{p+1}) T). 
\end{align*}
The analysis of the power series expansion of $\tilde{\mc P}(sT,c(sT)^{p+1}))^{1/s}$ reduces to analyzing the scaled effective generator \begin{align*}
    \frac{1}{sT} \sum_{q=1}^\infty \Phi_q(sT,c(sT)^{p+1}).
\end{align*}
To resolve the divergence issue of this BCH series as discussed in \sec{bch-truncate}, we consider the truncated scaled effective generator \begin{align}
\label{eq:def-effect-generator}
    \G_{(q_0)}(s) := \frac{1}{sT} \sum_{q=1}^{q_0} \Phi_q(sT,c(sT)^{p+1}),
\end{align}
where the truncation order $q_0$ is determined by \thm{general-bch-truncation} so that error induced by the truncation, $\|\exp({\G_{(q_0)}(s)T})-\tilde{\mc P}(sT,c(sT)^{p+1})^{1/s}\|$, is within the required error tolerance.
The Richardson extrapolation error is therefore determined by the power series of the truncated scaled effective generator $\G_{(q_0)}(s)$.

To analyze the power series of $\G_{(q_0)}(s)$, we first show that each $\Phi_q(t,\lambda)$ is a homogeneous polynomial of $t$ and $\lambda$ and bound the coefficients of the polynomial. For ease of presentation, we also define the generators without coefficients as
\begin{align}
\label{eq:def-K-tilde}
 	\tilde{\mathcal{K}}_j := \begin{cases}\ad_{H_{\gamma_v}} & \text{if } j = 2v-1\\ \mathcal{L}_v & \text{if } j = 2v \end{cases}
\end{align}
for $v \in [M]$.
\begin{lemma}[BCH Coefficient Bound]
\label{lem:bch-bound}
 	Define the sum of the commutator norms for the base generators $\tilde{\mc K}_j$ as
 	\begin{align}
 	\label{eq:alpha-q-q}
 	\alpha_{\comm}^{(q_1, q_2)}=\sum_{(j_1, \ldots, j_q)\in S_{q_1, q_2}}\big\|[\tilde{\mc K}_{j_1}, \tilde{\mc K}_{j_2},\ldots,\tilde{\mc K}_{j_q}]\big\|,
 	\end{align}
 	where $q = q_1+q_2$ and
 	\begin{align}
 	\label{eq:S-q-q}
 	 	S_{q_1, q_2}:=\Bigl\{(j_1, \ldots, j_q)\in [2M]^q\Bigm| &|\{k\in [q]\mid j_k \text{ is odd}\}|=q_1, |\{k\in [q]\mid j_k \text{ is even}\}|=q_2\Bigr\}. \nonumber
 	\end{align}
 	The $q$-th order term $\Phi_q(t,\lambda)$ in the BCH formula of $\tilde{\P}(t,\lambda)$ can be written as a homogeneous polynomial of $t$ and $\lambda$:
 	\begin{align}
 	\Phi_q(t,\lambda) = \sum_{\substack{q_1, q_2 \ge 0\\q_1+q_2 = q}}\Phi_{q_1, q_2} t^{q_1}\lambda^{q_2},
 	\end{align}
 	such that the coefficients $\Phi_{q_1, q_2}$ are bounded by
 	\begin{align}
 	\|\Phi_{q_1, q_2}\|\le \frac{a_{\max}^{q_1}}{q^2}\alpha_{\comm}^{(q_1, q_2)}.
 	\end{align}
\end{lemma}
\begin{proof}
 	Substituting Eq.~\eq{def-K} into the expression of the $q$-th order BCH term $\Phi_q(t,\lambda)$ (Eq.~\eq{bch}) yields 
 	\begin{align}
 	 	\Phi_q(t,\lambda)
 	 	=~ \begin{multlined}[t]
 	 	 	\sum_{\substack{p_v,p'_v\ge 0\\ p_1+p_1'+\cdots+p_M+p_M'=q }}\bigg(\prod_{v=1}^M \frac{\lambda^{p_v}(-\i ta_v)^{p'_v}}{p_v!p'_v!} \bigg) \\ 
 	 	\cdot \sum_{\sigma\in S_q} c_{\sigma, q}R_{\sigma}\big\{[\underbrace{\mc L_M, \ldots,\mc L_M}_{p_M}, \underbrace{\A_{\gamma_M}, \ldots, \mc A_{\gamma_M}}_{p_M'}, \ldots, \underbrace{ \mc L_1,  \ldots, \mc L_1}_{p_1}, \underbrace{\A_{\gamma_1},  \ldots, \mc A_{\gamma_1}}_{p_1'}]\big\}.\label{eq:Z_q}\end{multlined}
 	\end{align}
 	
 	Note that any summand in Eq.~\eq{Z_q} containing $[\tilde{\mc K}_{j_1}, \tilde{\mc K}_{j_2}, \ldots, \tilde{\mc K}_{j_q}]$ for $(j_1, \ldots, j_q)\in [2M]^q$ is a monomial of $\lambda$ and $t$. The degree of $t$ is $|\{k\in [q]\mid j_k \text{ is odd}\}|=: q_1$ and the degree of $\lambda$ is $|\{k\in [q]\mid j_k \text{ is even}\}|=: q_2$. We now bound the coefficient of this monomial. The commutator $[\tilde{\mc K}_{j_1}, \tilde{\mc K}_{j_2}, \ldots, \tilde{\mc K}_{j_q}]$ can only originate from one summand in the sum over $p_v, p'_v$ in Eq.~\eq{Z_q} because the operator counts $(p_1, p_1', \ldots, p_M, p_M')$ are uniquely determined by $(j_1, \ldots, j_q)$. The number of permutations $\sigma\in S_q$ in the second sum in Eq.~\eq{Z_q} that preserve the commutator is
 	\begin{align*}
 	 	\prod_{v=1}^M p_v!p'_v!.
 	\end{align*}
 	Thus, the coefficient of the monomial $t^{q_1}\lambda^{q_2}[\tilde{\mc K}_{j_1}, \ldots, \tilde{\mc K}_{j_q}]$ can be upper bounded by
 	\begin{align*}
 	 	\bigg\|\frac{1}{q^2}\prod_{v=1}^M \frac{(-\i a_v)^{p'_v}}{p_v!p'_v!}\prod_{v=1}^M p_v!p'_v!\big\|[\tilde{\mc K}_{j_1}, \tilde{\mc K}_{j_2}, \ldots, \tilde{\mc K}_{j_q}]\big\|\bigg\| \le \frac{a_{\max}^{q_1}}{q^2}  \big\|[\tilde{\mc K}_{j_1}, \tilde{\mc K}_{j_2}, \ldots, \tilde{\mc K}_{j_q}]\big\|,
 	\end{align*}
 	as $|c_{\sigma, q}|\le 1/q^2$. Then, we regroup the sum in Eq.~\eq{Z_q} according to $(j_1, \ldots, j_q)\in S_{q_1, q_2}$ (which have the same degree of $t$ and $\lambda$). The resulting polynomial
 	\begin{align*}
 	 	\Phi_q(t,\lambda) = \sum_{\substack{q_1, q_2 \ge 0\\q_1+q_2 = q}}\Phi_{q_1, q_2} t^{q_1}\lambda^{q_2}
 	\end{align*}
 	satisfies
 	\begin{align*}
 	 	\|\Phi_{q_1, q_2}\|\le  \frac{a_{\max}^{q_1}}{q^2}\sum_{(j_1, \ldots, j_q)\in S_{q_1, q_2}} \big\|[\tilde{\mc K}_{j_1}, \tilde{\mc K}_{j_2}, \ldots, \tilde{\mc K}_{j_q}]\big\| = \frac{a_{\max}^{q_1}}{q^2}\alpha_{\comm}^{(q_1, q_2)}.
 	\end{align*}
\end{proof}


By \lem{bch-bound}, $\G_{(q_0)}(s)$ is a power series in $s$ and we denote the $q$-th order coeffcient by $\E_{(q_0),q+1}$ such that \begin{align}
\label{eq:def-effect-generator}
    \G_{(q_0)}(s)=\sum_{q=0}^{\infty} \E_{(q_0),q+1} (sT)^q. 
\end{align}
We bound the norm of $\E_{(q_0),q}$ in the following lemma. 
\begin{lemma}
    \label{lem:bound-E}
    Define a constant
    \begin{align}\label{eq:def-mu}
        \mu_{\mathrm{comm},q_0} = \max_{q' \ge p+1} \Bigg(\sum_{\substack{1\le q_1+ q_2\le q_0 \\
        q_1+(p+1)q_2 = q'}}\frac{a_{\max}^{q_1}c^{q_2}}{(q_1+q_2)^2} \alpha_{\rm comm}^{(q_1, q_2)}\Bigg)^{1/q'}.
    \end{align}
    Then we have \begin{align}
        \G_{(q_0)}(s) = -\i \ad_H + \sum_{q=p}^{\infty} \E_{(q_0),q+1} (sT)^q \label{eq:def-E},
    \end{align} where $\E_{(q_0),q}$ can be bounded by \begin{align}
        \|\E_{(q_0),q}\| \le  \mu_{\mathrm{comm},q_0}^q.
    \end{align}
\end{lemma}
\begin{proof}
    By the order condition of the product formula (see \citep[Lemma 2]{watson2025exponentially}), we have \begin{align*}
        \frac{1}{sT}\sum_{q=1}^{\infty} \Phi_{q}(sT,0) = -\i \ad_H + O((sT)^p),
    \end{align*}
    which implies \begin{align*}
        \Phi_{q_1, 0} = 0,
    \end{align*}
    for all $1 <q_1 \le p$, and \begin{align*}
        \Phi_{1,0} = - \i \ad_H sT.
    \end{align*} 
    Therefore, for any $c>0$, we have \begin{align*}
        \frac{1}{sT} \Phi_{q}(sT,c(sT)^{p+1})&= \frac{1}{sT}\Big(\sum_{q_1=1}^{q_0}\Phi_{q_1, 0}(sT)^{q_1} + \sum_{q = 1}^{q_0}\sum_{\substack{q_1\ge 0,q_2\ge 1\\ q_1+q_2 = q}} \Phi_{q_1, q_2} (sT)^{q_1}(c(sT)^{p+1})^{q_2}\Big) \\
        & = -\i \ad_H + \sum_{q_1 = p+1}^{q_0} \Phi_{q_1, 0} (sT)^{q_1-1} + O((sT)^{p}) \\
        & = -\i \ad_H + \sum_{q = p}^{\infty} \E_{(q_0),q+1} (sT)^{q}.
    \end{align*}
    For any $q\ge p+1$, we have \begin{align*}
        \|\E_{(q_0),q}\| &= \bigg\|\sum_{q'=1}^{q_0}\sum_{\substack{q_1, q_2\ge 0,~ q_1+q_2=q' \\
        q_1+(p+1)q_2 = q}} \Omega_{q_1,q_2} c^{q_2} \bigg\|\\
        &\le \sum_{\substack{1\le q_1+ q_2\le q_0 \\
        q_1+(p+1)q_2 = q}}\frac{a_{\max}^{q_1}c^{q_2}}{(q_1+q_2)^2} \alpha_{\rm comm}^{(q_1, q_2)} \\
        &\le  \mu_{\mathrm{comm},q_0}^q.
    \end{align*}
    %For $q=p+1$, we have \begin{align*}
    %    \|\E_{p+1}\| &\le c\Big\|\sum_{v=1}^M \L_v\Big\|+\|\Omega_{p+1,0}\| \le c\alpha_{\rm comm}^{(0,1)}+\frac{a_{\max}^{p+1}}{(p+1)^2}\alpha_{\rm comm}^{(p+1,0)}\le \mu_{\mathrm{comm},q_0}^{p+1.}
    %\end{align*}
\end{proof}
\begin{lemma}
\label{lem:expansion-exp-G(s)}
    The evolution $e^{\G_{(q_0)}(s) T}$ admits series expansion in $s$ given by \begin{align*}
        e^{\G_{(q_0)}(s) T}=e^{-\i \ad_H T}+ \sum_{q=p}^{\infty}\tilde{\E}_{(q_0),q} s^q,
    \end{align*}
    where $\tilde{\E}_{(q_0),q}$ are superoperators bounded by \begin{align*}
       \|\tilde{\E}_{(q_0),q}\|\le   e(2\mu_{\mathrm{comm},q_0}T)^{q(1+1/p)},
    \end{align*}
    if $2\mu_{\mathrm{comm},q_0}T\ge 1$, and bounded by \begin{align*}
        \|\tilde{\E}_{(q_0),q}\|\le e(2\mu_{\mathrm{comm},q_0}T)^{q},
    \end{align*}
    if $2\mu_{\mathrm{comm},q_0}T\le 1$.
\end{lemma}
\begin{proof}
By the Dyson series expansion in the interaction picture, we have 
\begin{align*}
    &e^{\G_{(q_0)}(s) T}-e^{-\i \ad_H T} \\
    = \ & e^{-\i \ad_H T}  \sum_{j=1}^{\infty} \int_{0}^T \d \tau_1 \int_{0}^{\tau_1}\d \tau_2 \cdots \int_{0}^{\tau_{j-1}}\d \tau_j \prod_{j'=j}^{1}\bigg(e^{\i\ad_H \tau_{j'}}\bigg(\sum_{q=p}^{\infty}\E_{(q_0),q+1}s^qT^q\bigg)e^{-\i \ad_H \tau_{j'}}\bigg) \\
    = \ &  e^{-\i \ad_H T}\sum_{q=p}^{\infty}s^{q} \sum_{j=1}^{\infty} \sum_{\substack{q_1,\ldots,q_j \ge p \\ q_1+\cdots+q_j=q}} \int_{0}^T \d \tau_1 \int_{0}^{\tau_1}\d \tau_2 \cdots \int_{0}^{\tau_{j-1}}\d \tau_j \prod_{j'=j}^{1}\bigg(e^{\i\ad_H \tau_{j'}}\E_{(q_0),q_{j'}+1}e^{-\i \ad_H \tau_{j'}}T^{q_{j'}}\bigg)\\
    =:\,& \sum_{q=p}^{\infty}\tilde{\E}_{(q_0),q} s^q   
\end{align*}
The superoperator $\tilde{\E}_{(q_0),q}$ can be bounded by \begin{align*}
    \big\|\tilde{\E}_{(q_0),q}\big\|&\le \sum_{j=1}^{\infty} \sum_{\substack{q_1,\ldots,q_j \ge p \\ q_1+\cdots+q_j=q}} \int_{0}^T \d \tau_1 \int_{0}^{\tau_1}\d \tau_2 \cdots \int_{0}^{\tau_{j-1}}\d \tau_j \prod_{j'=1}^{j}\big(\big\|\E_{(q_0),q_{j'}+1}\big\|T^{q_{j'}}\big) \\
    &\le  \sum_{j=1}^{\infty} \sum_{\substack{q_1,\ldots,q_j \ge p \\ q_1+\cdots+q_j=q}}\frac{T^{q+j}}{j!}\prod_{j'=1}^{j}\mu_{\mathrm{comm},q_0}^{q_{j'}+1} \\
    &=\sum_{j=1}^{\lfloor q/p\rfloor }\frac{(\mu_{\mathrm{comm},q_0}T)^{q+j}}{j!}\Big(\sum_{\substack{q_1,\ldots,q_j \ge p \\ q_1+\cdots+q_j=q}} 1\Big).
\end{align*}
The second sum can be bounded by \begin{align*}
    \sum_{\substack{q_1,\ldots,q_j \ge p \\ q_1+\cdots+q_j=q}} 1 &\le \sum_{\substack{ q_1,\ldots,q_j \ge 0 \\ q_1+\cdots+q_j=q}} 1=\binom{q+j-1}{j-1} \le 2^{q+j}.
\end{align*}
Therefore, for $2\mu_{\mathrm{comm},q_0}T\ge 1$, we have \begin{align*}
    \big\|\tilde{\E}_{(q_0),q}\big\|\le \sum_{j=1}^{\lfloor q/p\rfloor }\frac{(2\mu_{\mathrm{comm},q_0}T)^{q+j}}{j!} 
    \le (2\mu_{\mathrm{comm},q_0}T)^{q+q/p}\sum_{j=1}^{\infty}\frac{1}{j!}\le e(2\mu_{\mathrm{comm},q_0}T)^{q(1+1/p)} ,
\end{align*}
and for $2\mu_{\mathrm{comm},q_0}T\le 1$, we have \begin{align*}
    \big\|\tilde{\E}_{(q_0),q}\big\| \le (2\mu_{\mathrm{comm},q_0}T)^{q}\sum_{j=0}^{\infty}\frac{ (2\mu_{\mathrm{comm},q_0}T)^{j}}{j!}\le (2\mu_{\mathrm{comm},q_0}T)^{q} e^{2\mu_{\mathrm{comm},q_0}T}\le e(2\mu_{\mathrm{comm},q_0}T)^{q}.
\end{align*}
\end{proof}
\xz{Also assume that the noise channel after each layer can be implemented efficiently. }
\sz{Here noise channel arises naturally, not like in Lindbladian simulation as modeling objective to be simulated/implemented.}

We now combine the BCH truncation bound, the effective-generator expansion, and Richardson extrapolation to obtain a general resource bound for the joint mitigation protocol.
\begin{theorem}
\label{thm:general-application}
    Let $H = \sum_{\gamma = 1}^{\Gamma} H_{\gamma}$ be the target Hamiltonian. Suppose we approximate its evolution by implementing a $p$-th order product formula with $M$ gate layers:
    \begin{align*}
        P(t) = \prod_{v=1}^M e^{-\i t a_v H_{\gamma_v}},
    \end{align*}
    and that the implementation is subject to physical noise described by $e^{\lambda \L_v}$ after the $v$-th gate layer, with a tunable noise strength $\lambda \ge \lambda_{\min}$.
    
    Let $c>0$ and $q_0\in \mb Z_+$. Let $\bar{\alpha}_{\mathrm{comm},q_0}(s)$ be the doubly right-nested commutator bound (Eq.~\eq{def-beta-comm-sum}) evaluated for the generators:
    \begin{align*}
        \mathcal{K}_j= \begin{cases}-\i sT a_{v}\ad_{H_{\gamma_v}} & \text{if } j = 2v-1\\ c(sT)^{p+1}\mathcal{L}_v & \text{if } j = 2v \end{cases}.
    \end{align*}
    Let $\mu_{\mathrm{comm},q_0}$ be the parameter defined in Eq.~\eq{def-mu}. 
    
    There exists an algorithm that, given any observable $O$ and initial state $\rho_0$, estimates $\operatorname{Tr}[O e^{-\i \ad_H T}\rho_0]$ to within an additive error of $\epsilon\|O\|$ using a total of \begin{align*}
        O\Big( \max\{\log^3(1/\eps), \log^2(1/\eps)s_0^{-1}\}\frac{(\log\log (1/\eps))^3}{\eps^2}\Big)  = \widetilde{O}\Big(\frac{1}{s_0 \eps^2}\Big)
    \end{align*}Trotter steps. This involves $\widetilde{O}(1/\epsilon^2)$ independent circuit runs, each comprising at most 
    \begin{align*}
        O( \max\{\log^3(1/\eps), \log^2(1/\eps)s_0^{-1}\}) = \widetilde{O}\left( \frac{1}{s_0} \right)
    \end{align*} Trotter steps,
    where $s_0$ is the largest step size satisfying
    \begin{align} \label{eq:s0-general}
       s_0 (2\mu_{\mathrm{comm},q_0}T)^{1+1/p} \le 1/e \quad \text{and} \quad
       2e \bar{\alpha}_{\mathrm{comm},q_0}(s_0) \le s_0 \frac{\epsilon}{4eC\log m},
    \end{align}
    with $C$ being the Richardson constant and $m = O(\log(1/\eps))$ satisfying $2e^{1-m} \le \eps /(4C\log m)$.
    
    This result holds provided that the intrinsic noise strength of the device, $\lambda_{\min}$, satisfies the feasibility condition $\lambda_{\min} = \widetilde{O}(c (s_0 T)^{p+1})$.
\end{theorem}
\begin{proof}
    Recall that the function for extrapolation is \begin{align*}
        f(s)&=\operatorname{Tr}[O(\tilde{\mc P}(sT, c(sT)^{p+1}))^{1/s}\rho_0],
    \end{align*}
    for $s>0$ and $1/s\in \mb Z$.
    Here we consider the scenario where $2\mu_{\mathrm{comm},q_0} T \ge 1$. By \lem{expansion-exp-G(s)}, the evolution $e^{\G_{(q_0)}(s) T}$ admits series expansion in $s$ given by \begin{align*}
        e^{\G_{(q_0)}(s) T}=e^{-\i \ad_H T}+ \sum_{q=p}^{\infty}\tilde{\E}_{(q_0),q} s^q,
    \end{align*}
    such that \begin{align*}
        \|\E_{(q_0), q}\|\le e (2\mu_{\mathrm{comm},q_0}T)^{q(1+1/p)}.
    \end{align*}    
    Then the function $f(s)$ can be decomposed as \begin{align*}
        f(s)&= \operatorname{Tr}[Oe^{\G_{(q_0)}(s)T} \rho_0]+\operatorname{Tr}[O\bigl((\tilde{\mc P}(sT, c(sT)^{p+1}))^{1/s}-e^{\G_{(q_0)}(s)T}\bigr)\rho_0] \\
        &= \underbrace{\operatorname{Tr}[Oe^{-\i \ad_H T}\rho_0]+\sum_{q=p}^{m-1}\operatorname{Tr}[O \E_{(q_0), q}\rho_0] s^q}_{P_m(s)}+\underbrace{\sum_{q=m}^{\infty}\operatorname{Tr}[O\E_{(q_0), q}\rho_0] + \operatorname{Tr}[O(\tilde{\mc P}(sT, c(sT)^{p+1})-e^{\G_{(q_0)}(s)T})\rho_0]}_{R_m(s)},
    \end{align*}
    where $P_m$ is a degree-$(m-1)$ polynomial in $s$. Let $s_1, \ldots, s_m$ and $b_1, \ldots, b_m$ be the extrapolation points and coefficients defined in \lem{richardson-extrapolation} with $t=1$ and $s = s_0$ to be determined later, and then we have $1\le\|\mathbf{b}\|\le C\log m$, $1/s_j\in \mathbb{Z}$, and $1/s_j = r_j = O(\max\{m^3, m^2/s\}/j^2)$. We can set \begin{align*}
        m = O(\log(1/\eps)),
    \end{align*}
    such that \begin{align}
    \label{eq:cond-m}
        2e^{1-m} \le \frac{\eps}{4C\log m}.
    \end{align}
    
   
    For any $s$ satisfying \begin{align}
    \label{eq:cond-s}
        e\bar\alpha_{\mathrm{comm},q_0}(s)\le s\frac{\eps}{4eC \log m} \le s \quad s(2\mu_{\mathrm{comm},q_0}T)^{1+1/p} \le 1/e,
    \end{align} the remainder $R_m(s)$ is bounded as \begin{align*}
        &|R_m(s)| \\
        \le ~&\sum_{q=m}^{\infty}\|O\|\|\tilde{\E}_{(q_0),q}\|\|\rho_0\|_1 s^q + \|O\|\Big\|(\tilde{\mc P}(sT, c(sT)^{p+1}))^{1/s}-(e^{\G_{(q_0)}(s)sT})^{1/s}\Big\|\|\rho_0\|_1 \\
        \le ~& \sum_{q=m}^{\infty} e(s(2\mu_{\mathrm{comm},q_0} T)^{1+1/p})^{q}\|O\|+\frac{(\max\{1, \|e^{\G_{(q_0)}(s) sT}\|\})^{1/s}}{s}\Bigl\|\tilde{\mc P}(sT, c(sT)^{p+1})-e^{\G_{(q_0)}(s)sT}\Bigr\|\|O\|\\
        \le ~&2e^{1-m} \|O\|+ \frac{(1+e\alpha_{\mathrm{comm},q_0}(s))^{1/s}}{s}e\alpha_{\mathrm{comm},q_0}(s)\|O\| \\
        \le ~& \frac{\eps}{4C\log m} \|O\|+(1+s)^{1/s}\frac{\eps}{4eC\log m}\|O\| \le \frac{\eps}{2C\log m} \|O\|
    \end{align*}
    where the third line follows from \lem{expansion-exp-G(s)} and \lem{power-bound}, the fourth line follows from 
    \thm{general-bch-truncation}, and the fifth line follows from Eq.~\eq{cond-m} and Eq.~\eq{cond-s}.
    
    Now, we bound the error of the extrapolation function \begin{align*}
        F^{(m)} = \sum_{j=1}^m b_j f(s_j).
    \end{align*}
    Set \begin{align*}
        f(0) :=\lim_{s\to 0} f(s) = \operatorname{Tr}[Oe^{-\i \ad_H T}\rho_0].
    \end{align*} Let $s_0$ be the largest number smaller than $e^{-1}(2\mu_{\mathrm{comm},q_0}T)^{-1-1/p}$ and satisfying \begin{align*}
        e\bar\alpha_{\mathrm{comm},q_0}(s_0)\le s_0\frac{\eps}{4eC \log m}
    \end{align*}
    Observe that $\bar\alpha_{\mathrm{comm}, q_0}(s)$ consists of sums of norms of nested commutators involving $-\i sT a_v \mc{A}_{\gamma_v}$ and $c(sT)^{p+1} \mc L_{v}$ with grade $q \ge q_0+1$. Consequently, for $s \ge 0$, it is a polynomial in $s$ with non-negative coefficients and a lowest degree of $q_0+1$. This implies that the function $\bar\alpha_{\mathrm{comm}, q_0}(s)/s$ is monotonically increasing for $s \ge 0$. 
    By \lem{richardson-extrapolation}, the extrapolation nodes satisfy $s_j \le s_0$. Therefore, as the condition in Eq.~\eq{cond-s} holds for the base step size $s_0$, it is satisfied for all $s_j$.
    Then we have \begin{align*}
        |F^{(m)}(s)-f(0)|&\le \|\mathbf{b}\|_1 \max_j|R_m(s_j)| \\
        &\le C\log m\frac{\eps}{2C\log m}\|O\|\\
        &\le \frac{\eps }{2}\|O\|.
    \end{align*}
    To estimate $f(s_j)$ with error $\eps\|O\|/(2\|\mathbf{b}\|_1)$ with success probability $1-1/(3m)$, we need to repeatedly run $(\tilde{\mc P}(s_jT, c(s_jT)^{p+1}))^{1/s_j}$ $O(\|\mathbf{b}\|_1^2/\eps^2\log m)$ times.  Each application of $(\tilde{\mc P}(s_jT, c(s_jT)^{p+1}))^{1/s_j}$ requires $1/s_j$ Trotter steps. Hence, the maximum number of Trotter steps per coherent run is \begin{align*}
        O(\max_j r_j) & = O( \max\{m^3, m^2s_0^{-1}\})\\
        &=O( \max\{\log^3(1/\eps), \log^2(1/\eps)s_0^{-1}\}) \\
        &= \widetilde{O}\Big(\frac{1}{s_0}\Big),
    \end{align*} and the total number of Trotter steps is \begin{align*}
        O\Big(\sum_{j=1}^m r_j\frac{\|\mathbf{b}\|_1^2}{\eps^2}\log m\Big) & = O\Big(\sum_j^{m} \max\{m^3, m^2s_0^{-1}\}\frac{1}{j^2}\frac{\log^3 m}{\eps^2}\Big)\\
        &=O\Big( \max\{\log^3(1/\eps), \log^2(1/\eps)s_0^{-1}\}\frac{(\log\log (1/\eps))^3}{\eps^2}\Big) \\
        &= \widetilde{O}\Big(\frac{1}{s_0\eps^2}\Big)
    \end{align*}
    The smallest noise strength is \begin{align*}
        \min_{j}c(s_jT)^{p+1} &= \Omega\Big(\min_{j}c\Big(\frac{j^2T}{\max\{m^3, m^2s_0^{-1}\}}\Big)^{p+1}\Big) \\
        &=\Omega\Big(c\min\Big\{\frac{T^{p+1}}{\log^{3p+3}(1/\eps)}, \frac{ T^{p+1}s_0^{p+1}}{\log^{2p+2}(1/\eps)}\Big\}\Big)\\
        & = \widetilde{\Omega}(c (s_0 T)^{p+1}).
    \end{align*}
\end{proof}
\xz{TBD: explaining three tunable parameters: $c$ balances Trotter error and physical error, $q_0$ balances BCH truncation error and extrapolation error.}
The algorithm involves three tunable parameters that control different aspects of the error budget:
\begin{itemize}
    \item The coupling constant $c$ in $\lambda(s)=c(sT)^{p+1}$ determines the relative weight of the physical noise error versus the Trotter error in the leading-order term. In \sec{local}, we optimize $c$ analytically (Eq.~\eqref{eq:def-c}) to balance the two sources so that neither dominates the extrapolation remainder.
    \item The BCH truncation order $q_0$ controls the error from approximating the full BCH series by its truncation (see \thm{general-bch-truncation}). A larger $q_0$ tightens the doubly-nested commutator bound $\bar{\alpha}_{\mathrm{comm},q_0}$ at the cost of a slightly larger commutator parameter $\mu_{\mathrm{comm},q_0}$ in the effective generator expansion. We choose $q_0 = O(\log(N/\eps))$, which suffices to suppress the truncation error below the target tolerance.
    \item The number of extrapolation nodes $m$ (equivalently, the minimum step size $s_0$) governs the Richardson extrapolation accuracy. By \lem{richardson-extrapolation}, the extrapolation error decays exponentially in $m$, so $m = O(\log(1/\eps))$ nodes are sufficient.
\end{itemize}

\sz{\paragraph{Variance-sensitive sampling refinement.}
The shot-complexity estimate above uses a worst-case bounded-outcome concentration argument. A sharper statement keeps the single-shot variances of the extrapolation data points. We record this refinement separately because it clarifies how a one-dimensional path can have a smaller statistical prefactor than a two-dimensional grid, even though both retain the same leading $1/\eps^2$ dependence in the worst case.
}
\begin{lemma}[Bernstein bound with shot allocation]
\label{lem:bernstein-shot-allocation}
Let
\begin{align*}
    \widehat F=\sum_{i=1}^{K}a_i\widehat \mu_i,
    \qquad
    \widehat\mu_i=\frac{1}{n_i}\sum_{\ell=1}^{n_i}X_{i,\ell},
\end{align*}
where the random variables $X_{i,\ell}$ are independent, $\mb E X_{i,\ell}=\mu_i$, $\mathrm{Var}(X_{i,\ell})=\sigma_i^2$, and $X_{i,\ell}\in[-\|O\|,\|O\|]$. Then
\begin{align}
\label{eq:bernstein-linear-estimator}
    \Pr\{|\widehat F-\mb E\widehat F|\ge \eta\}
    \le
    2\exp\left(
    -\frac{\eta^2}{2V+\frac{4}{3}\|O\|\eta B}
    \right),
\end{align}
where
\begin{align*}
    V=\sum_{i=1}^{K}\frac{a_i^2\sigma_i^2}{n_i},
    \qquad
    B=\max_i\frac{|a_i|}{n_i}.
\end{align*}
If every data point receives the same number of shots, the total number of shots sufficient for statistical error at most $\eta$ with failure probability at most $\delta$ is
\begin{align}
\label{eq:equal-shot-bound}
    N_{\mathrm{eq}}
    =
    O\left[
    \left(
    \frac{K\sum_{i=1}^K a_i^2\sigma_i^2}{\eta^2}
    +
    \frac{\|O\|K\max_i|a_i|}{\eta}
    \right)
    \log\frac{1}{\delta}
    \right].
\end{align}
Moreover, by allocating shots according to the Cauchy--Schwarz, or Neyman, rule for the variance term, with an additional floor proportional to $|a_i|$ to control the range term, one obtains
\begin{align}
\label{eq:optimized-shot-bound}
    N_{\mathrm{opt}}
    =
    O\left[
    \left(
    \frac{\left(\sum_{i=1}^{K}|a_i|\sigma_i\right)^2}{\eta^2}
    +
    \frac{\|O\|\|\mathbf a\|_1}{\eta}
    \right)
    \log\frac{1}{\delta}
    \right].
\end{align}
\end{lemma}
\begin{proof}
Eq.~\eqref{eq:bernstein-linear-estimator} follows by applying Bernstein's inequality to the centered variables $a_i(X_{i,\ell}-\mu_i)/n_i$. Their total variance is $V$, and each summand has absolute value at most $2\|O\||a_i|/n_i$.

For equal shots, $n_i=N/K$, and therefore $V=K\sum_i a_i^2\sigma_i^2/N$ and $B=K\max_i|a_i|/N$, which gives Eq.~\eqref{eq:equal-shot-bound}. For the optimized allocation, let $S=\sum_i |a_i|\sigma_i$. Allocate one part of the shot budget with $n_i^{(v)}\propto |a_i|\sigma_i$ and one part with $n_i^{(r)}\propto |a_i|$. The first part gives $V\le S^2/N_v$ by Cauchy--Schwarz, while the second gives $B\le \|\mathbf a\|_1/N_r$. Choosing $N_v=O(S^2\eta^{-2}\log(1/\delta))$ and $N_r=O(\|O\|\|\mathbf a\|_1\eta^{-1}\log(1/\delta))$ yields Eq.~\eqref{eq:optimized-shot-bound}.
\end{proof}

For the one-dimensional joint extrapolation, take $a_j=b_j$ and
\begin{align*}
    \sigma_{j,\mathrm{path}}^2
    :=
    \mathrm{Var}\!\left(X_{s_j,\lambda(s_j)}\right).
\end{align*}
The equal-shot and optimized shot complexities become
\begin{align}
\label{eq:1d-equal-shot}
    N_{\mathrm{1D,eq}}
    =
    O\left[
    \left(
    \frac{m\sum_j b_j^2\sigma_{j,\mathrm{path}}^2}{\eta^2}
    +
    \frac{\|O\|m\max_j|b_j|}{\eta}
    \right)
    \log\frac{1}{\delta}
    \right],
\end{align}
and
\begin{align}
\label{eq:1d-opt-shot}
    N_{\mathrm{1D,opt}}
    =
    O\left[
    \left(
    \frac{\left(\sum_j |b_j|\sigma_{j,\mathrm{path}}\right)^2}{\eta^2}
    +
    \frac{\|O\|\|\mathbf b\|_1}{\eta}
    \right)
    \log\frac{1}{\delta}
    \right].
\end{align}
For a tensor-product two-dimensional extrapolation with step-size coefficients $b_j$ and noise-direction coefficients $d_\ell$, the estimator has coefficients $a_{j,\ell}=b_jd_\ell$. Writing
\begin{align*}
    \sigma_{j,\ell}^2
    :=
    \mathrm{Var}\!\left(X_{s_j,\lambda_\ell}\right),
\end{align*}
we similarly get
\begin{align}
\label{eq:2d-equal-shot}
    N_{\mathrm{2D,eq}}
    =
    O\left[
    \left(
    \frac{m_sm_\lambda\sum_{j,\ell} b_j^2d_\ell^2\sigma_{j,\ell}^2}{\eta^2}
    +
    \frac{\|O\|m_sm_\lambda\max_{j,\ell}|b_jd_\ell|}{\eta}
    \right)
    \log\frac{1}{\delta}
    \right],
\end{align}
and
\begin{align}
\label{eq:2d-opt-shot}
    N_{\mathrm{2D,opt}}
    =
    O\left[
    \left(
    \frac{\left(\sum_{j,\ell}|b_jd_\ell|\sigma_{j,\ell}\right)^2}{\eta^2}
    +
    \frac{\|O\|\|\mathbf b\|_1\|\mathbf d\|_1}{\eta}
    \right)
    \log\frac{1}{\delta}
    \right].
\end{align}
Thus the Bernstein estimate does not by itself change the worst-case asymptotic scaling in $\eta$, but it replaces the range-only prefactor by the actual variance profile of the measured extrapolation data.

For a single Pauli observable, one shot gives $X\in\{\pm1\}$ and hence
\begin{align}
\label{eq:pauli-variance}
    \sigma^2=\mathrm{Var}(X)=1-\mu^2,
    \qquad
    \mu=\mb E X.
\end{align}
For the average magnetization $M_Z=N^{-1}\sum_{i=1}^{N}Z_i$, Eq.~\eqref{eq:pauli-variance} holds exactly if one uses a random-site Pauli estimator: choose $i$ uniformly at random, measure $Z_i$, and output $X=z_i\in\{\pm1\}$. Then $\mb E X=\langle M_Z\rangle$. If instead each shot measures the full bit string and outputs $Y=N^{-1}\sum_i z_i$, then
\begin{align*}
    \mathrm{Var}(Y)
    =
    \frac{1}{N^2}\sum_{i,k}\mathrm{Cov}(Z_i,Z_k),
\end{align*}
so correlations enter and Eq.~\eqref{eq:pauli-variance} is not generally an equality. Since $Y\in[-1,1]$, one still has $\mathrm{Var}(Y)\le 1-(\mb E Y)^2$, and the actual full-bit-string estimator can be treated through the empirical variance profile $\sigma_i^2$ in the bounds above.

In the random-site Pauli formulation, a sufficient condition for the one-dimensional path to have no larger Bernstein variance proxy than the two-dimensional grid is
\begin{align}
\label{eq:variance-proxy-condition-equal}
    \sum_j b_j^2\left(1-\mu_{j,\mathrm{path}}^2\right)
    \le
    \sum_{j,\ell}b_j^2d_\ell^2\left(1-\mu_{j,\ell}^2\right),
\end{align}
where $\mu_{j,\mathrm{path}}=\mb E X_{s_j,\lambda(s_j)}$ and $\mu_{j,\ell}=\mb E X_{s_j,\lambda_\ell}$. A stronger pointwise sufficient condition is
\begin{align}
\label{eq:variance-proxy-condition-pointwise}
    1-\mu_{j,\mathrm{path}}^2
    \le
    \sum_\ell d_\ell^2\left(1-\mu_{j,\ell}^2\right)
    \qquad \text{for every } j.
\end{align}
For the optimized allocation, the analogous sufficient condition is
\begin{align}
\label{eq:variance-proxy-condition-opt}
    \sum_j |b_j|\sqrt{1-\mu_{j,\mathrm{path}}^2}
    \le
    \sum_{j,\ell}|b_jd_\ell|\sqrt{1-\mu_{j,\ell}^2}.
\end{align}
These conditions formalize the intuition that the coupled path is statistically favorable when its measured points have smaller shot variance, while the two-dimensional grid pays the extra noise-direction coefficient weights.

\subsection{Applications to $k$-local Hamiltonians}
\label{sec:local}
Suppose that $\ad_H=\sum_{\gamma\in[\Gamma]}\A_\gamma$ and $\mc L_{\tot}:=\sum_{v\in [M]}\mc L_v$ are $k$-local Lindbladians on a lattice $\Lambda = [N]$. We further assume that $\ad_H$ is $g_H$-extensive and $\L_{\rm tot}$ is $g_L$-extensive.
Recall the unified representation of $\mc A_{\gamma}$ and $\mc L_v$ in Eq.~\eq{def-K} and Eq.~\eq{def-K-tilde}: \begin{align*}
    \mathcal{K}_j = \begin{cases}-\i ta_{v}\ad_{H_{\gamma_v}} & \text{if } j = 2v-1 \\ \lambda\mathcal{L}_v & \text{if } j = 2v \end{cases}, \quad \tilde{\mathcal{K}}_j := \begin{cases}\ad_{H_{\gamma_v}} & \text{if } j = 2v-1\\ \mathcal{L}_v & \text{if } j = 2v \end{cases} 
\end{align*}
for $v=1, 2,\ldots, M$.  

Now we substitute $t = sT, \lambda = c (sT)^{p+1}$ and let $\G_{(q_0)}(s)$ be the effective generator defined in Eq.~\eq{def-effect-generator}. By \lem{bound-E}, the coefficients of the series expansion of $e^{\G_{(q_0)}(s)T}$ can be bounded as powers of $\mu_{\mathrm{comm},q_0}$ defined in Eq.~\eq{def-mu}. We now provide an upper bound of $\mu_{\mathrm{comm},q_0}$. 
\begin{lemma}
\label{lem:bound-mu}
    The parameter $\mu_{\mathrm{comm},q_0}$ defined in Eq.~\eq{def-mu} satisfies \begin{align*}
        \mu_{\mathrm{comm},q_0} \le q_0 N^{1/(p+1)}(2a_{\max}k \Upsilon g_H+ (2ckg_L)^{1/(p+1)}).
    \end{align*}
\end{lemma}
\begin{proof}
    For the set $S_{q_1, q_2}$ defined in Eq.~\eq{alpha-q-q}, we decompose it according to the position of the odd and even indices.  For any index set $\mathcal{I} \subseteq [q_1+q_2]$, we define the subset $S_{\mathcal{I}}$ as
    \[
    S_{\mathcal{I}} = \Bigl\{ (j_1, \ldots, j_{q_1+q_2}) \in [2M]^{q_1+q_2} \Bigm| 
        j_k \text{ is odd for } k \in \mathcal{I}, \text{ and } 
        j_k \text{ is even for } k \in [q_1+q_2] \setminus \mathcal{I}
    \Bigr\}.
    \]
    We can then express $S_{q_1, q_2}$ as the disjoint union over the collection of all size-$q_1$ subsets of $[q_1+q_2]$:
    \[
    S_{q_1, q_2} = \bigsqcup_{\mathcal{I} \in \binom{[q_1+q_2]}{q_1}} S_{\mathcal{I}}.
    \]
    The commutator bounds are \begin{align*}
    \alpha_{\rm comm}^{(q_1,q_2)} &=  \sum_{(j_1, \ldots, j_{q_1+q_2})\in S_{q_1, q_2}}\big\|[\tilde{\mc K}_{j_1}, \tilde{\mc K}_{j_2},\ldots,\tilde{\mc K}_{j_{q_1+q_2}}]\big\|\\
    &= \sum_{\mc I\in \binom{[q_1+q_2]}{q_1}}\sum_{(j_1, \ldots, j_{q_1+q_2})\in S_{\mc I}}\big\|[\tilde{\mc K}_{j_1}, \tilde{\mc K}_{j_2},\ldots,\tilde{\mc K}_{j_{q_1+q_2}}]\big\|\\
    &\le \sum_{\mc I\in \binom{[q_1+q_2]}{q_1}} N(q_1+q_2-1)!(2k)^{q_1+q_2-1}(\Upsilon g_H)^{q_1}g_L^{q_2}\\
    &= \binom{q_1+q_2}{q_1}N(q_1+q_2-1)!(2k)^{q_1+q_2-1}(\Upsilon g_H)^{q_1}g_L^{q_2}\\
    &\le \binom{q_1+(p+1)q_2}{q_1}N(q_1+q_2)!(2k\Upsilon g_H)^{q_1}((2kg_L)^{1/(p+1)})^{q_2(p+1)},
    \end{align*}
    where the third line follows from \cor{commu-bound-sum}. Then we have  \begin{align*}
        \mu_{\mathrm{comm},q_0} &= \max_{q' \ge p+1} \Bigg(\sum_{\substack{1\le q_1+q_2\le q_0 \\
        q_1+(p+1)q_2 = q'}}\frac{a_{\max}^{q_1}c^{q_2}}{(q_1+q_2)^2} \alpha_{\rm comm}^{(q_1, q_2)}\Bigg)^{1/q'} \\
        &\le \max_{q' \ge p+1} \bigg(\sum_{\substack{1\le q_1+q_2\le q_0 \\
        q_1+q_2(p+1) = q'}}\binom{q_1+q_2(p+1)}{q_1}N(q_1+q_2)!(2a_{\max}k\Upsilon g_H)^{q_1}((2ckg_L)^{1/(p+1)})^{q_2(p+1)}\bigg)^{1/q'}\\
        &\le \max_{q' \ge p+1} \bigg(\sum_{q_1=1}^{q_0}\binom{q'}{q_1}Nq_0^{q'}(2a_{\max}k\Upsilon g_H)^{q_1}((2ckg_L)^{1/(p+1)})^{q'-q_1}\bigg)^{1/q'}\\
        &\le q_0N^{1/(p+1)}\max_{q'\ge p+1} \big((2a_{\max}k\Upsilon g_H+(2ckg_L)^{1/(p+1)})^{q'}\big)^{1/q'}\\
        &=q_0N^{1/(p+1)}(2a_{\max}k\Upsilon g_H+(2ckg_L)^{1/(p+1)}).
    \end{align*}
\end{proof}

Now we apply \thm{general-application} to $k$-local Hamiltonians. In the following theorem, we make the dependence on locality, extensiveness, and intrinsic noise explicit.  
\begin{theorem}
\label{thm:extrapolation-k-local}
    Suppose that the intrinsic noise strength of the device, $\lambda_{\min}$, satisfies
    \begin{align}
    \label{eq:cond-lambda}
        \lambda_{\min} = \widetilde{O}\left(\frac{1}{(Nk\Upsilon a_{\max}g_HT)^{1+1/p} kg_L}\right).
    \end{align} 
    There exists an algorithm that, given any observable $O$ and initial state $\rho_0$, estimates $\operatorname{Tr}[O e^{-\i \ad_H T}\rho_0]$ to within an additive error of $\epsilon\|O\|$ using a total of 
    \begin{align*}
        O\left( \max\left\{\log^3(1/\eps), \Xi \log^2(1/\eps)\right\} \frac{(\log\log (1/\eps))^3}{\eps^2} \right)
    \end{align*}
    Trotter steps, where we define the parameter-dependent factor
    \begin{align*}
        \Xi := N^{1/p}(k\Upsilon a_{\max} g_H T)^{1+1/p} \log^{1+1/p}(N/\eps).
    \end{align*}
    This involves $\widetilde{O}(1/\epsilon^2)$ independent circuit runs, each comprising at most
    \begin{align*}
        O\left( \max\left\{\log^3(1/\eps), \Xi \log^2(1/\eps)\right\} \right)
    \end{align*}
    Trotter steps.
\end{theorem}

\begin{proof}
    Note that $\mc K:= \sum_{j=1}^{2M} {\mc K}_j$ is $k$-local and $(\Upsilon sT a_{\max}g_H+c(sT)^{p+1} g_L)$-extensive.
    Following the proof of \lem{bch-approx}, the doubly right-nested commutator bound of $\mc K$, denoted $\bar{\alpha}_{\mathrm{comm}, q_0}(s)$, satisfies 
    \begin{align*}
        \bar{\alpha}_{\mathrm{comm}, q_0}(s) \le 2(4eq_0 k(\Upsilon sT a_{\max}g_H+c(sT)^{p+1} g_L))^{q_0+1}N,
    \end{align*}
    provided $4eq_0 k(\Upsilon sT a_{\max}g_H+c(sT)^{p+1} g_L) \le 1/2$. 
    Let $m=O(\log(1/\eps))$ be the number of extrapolation points. We choose the expansion order as $q_0 = \lceil \log(8e^2N\log m/\eps) \rceil = O(\log(N/\eps))$. 
    
    We set the step size $s_0$ as
    \begin{align}
    \label{eq:s0-def}
        s_0 &= \frac{1}{(2eq_0 N^{1/(p+1)}T (4ek \Upsilon a_{\max} g_H + (4ekcg_L )^{1/(p+1)}))^{1+1/p}} \\
        & = \Theta\left( \frac{1}{ \log^{1+1/p}(N/\eps)  N^{1/p}T^{1+1/p}(k\Upsilon a_{\max} g_H +(kcg_L)^{1/(p+1)})^{1+1/p} } \right).\nonumber
    \end{align}
    One can verify that this choice of $s_0$ satisfies the condition $4eq_0 k(\Upsilon s_0T a_{\max}g_H+c(s_0T)^{p+1} g_L) \le 1/2$.
    Furthermore, we have
    \begin{align*}
         \frac{\bar{\alpha}_{\mathrm{comm}, q_0}(s_0)}{s_0} & \le 2\max\{(s_0^{p/(p+1)}4eq_0k\Upsilon T a_{\max} g_H)^{q_0+1}, (4eq_0kcg_L s_0 ^pT^{p+1})^{q_0+1}\}N \\
         &\le 2e^{-q_0-1}N \le \frac{\eps}{8e^2C \log m}.
    \end{align*}
    Also, $s_0$ satisfies the condition related to $\mu_{\mathrm{comm}, q_0}$:
    \begin{align*}
        s_0 &\le \frac{1}{e(2q_0 N^{1/(p+1)} T(2 a_{\max} k\Upsilon g_H + (2ckg_L)^{1/(p+1)}))^{1+1/p}} \le \frac{1}{e(2\mu_{\mathrm{comm}, q_0}T)^{1+1/p}}. 
    \end{align*}
    
    To optimize the constraint on $\lambda_{\min}$ without increasing the asymptotic complexity, we set
    \begin{align}
    \label{eq:def-c}
        c = \frac{(k\Upsilon a_{\max}g_H)^{p+1}}{kg_L},
    \end{align}
    which implies $(kcg_L)^{1/(p+1)} = k\Upsilon a_{\max}g_H$. Substituting this into Eq.~\eqref{eq:s0-def}, the inverse step size scales as:
    \begin{align*}
        \frac{1}{s_0} = O\left( N^{1/p}(k\Upsilon a_{\max} g_H T)^{1+1/p} \log^{1+1/p}(N/\eps) \right) = O(\Xi).
    \end{align*}

    Now, applying \thm{general-application}, the algorithm requires a total number of Trotter steps of:
    \begin{align*}
         O\left( \max\left\{\log^3(1/\eps), \Xi \log^2(1/\eps)\right\} \frac{(\log\log (1/\eps))^3}{\eps^2} \right).
    \end{align*}
    The maximum number of Trotter steps per run is
    \begin{align*}
  O\left( \max\left\{\log^3(1/\eps), \Xi \log^2(1/\eps)\right\} \right).
    \end{align*}
    This result holds provided that the intrinsic noise strength satisfies:
    \begin{align*}
        \lambda_{\min} = \widetilde{O}(c(s_0T)^{p+1}) = \widetilde{O}\Big(\frac{1}{ (Nk\Upsilon a_{\max}g_HT)^{1+1/p} kg_L }\Big).
    \end{align*}
\end{proof}
\xz{This reaches the limit of error mitigation under Pauli noise \cite{wang2021noise, cai2021multi}.}
Refs.~\cite{wang2021noise, cai2021multi} established that under Pauli noise, any unbiased QEM protocol requires a sampling overhead of at least $\Omega(1/\epsilon^2)$ to achieve additive error $\epsilon$, and that the maximum circuit depth tolerable at a given noise rate is fundamentally limited. Our joint mitigation protocol achieves the optimal $\widetilde{O}(1/\epsilon^2)$ sampling overhead, and the feasibility condition on $\lambda_{\min}$ in Eq.~(\ref{eq:cond-lambda}) matches the fundamental noise threshold identified in these works. Thus, under the mild-noise assumption, our algorithm operates at the information-theoretic limit of what QEM can achieve without resorting to full quantum error correction.
In the noiseless limit ($g_L = 0$), \thm{extrapolation-k-local} provides a resource analysis of the Trotter error mitigation algorithm proposed by \citet{watson2025exponentially} for $k$-local lattice Hamiltonians. This result parallels that of \citet{mizuta2025commutator}, which analyzed the complexity of multi-product formulas without physical noise using a different proof technique.

A crucial implication of \thm{extrapolation-k-local} is that, even when the device is subject to local noise satisfying Eq.~\eq{cond-lambda}, the noiseless extrapolation algorithm can be applied with only a constant complexity overhead by properly coupling the Trotter step size and the noise strength in the extrapolation. This implies that physical noise mitigation is achieved at virtually no asymptotic cost. Specifically, our choice of $c$ in Eq.~\eq{def-c} merely rescales the step size $s_0$ by a factor of $2^{-1-1/p}$. Consequently, the total complexity, which is proportional to $s_0^{-1}$, increases only by a constant factor of $2^{1+1/p}$.

\xz{I'm thinking whether to explicitly compare our joint mitigation algorithm with the nested mitigation one. Since the asymptotic complexity advantage of our algorithm is only $\log(1/\eps)(\log\log(1/\eps))^2$. If not, we can instead claim that our algorithm can mitigate the physical error almost for free while doing Trotter error mitigation. }


\subsection{Applications to geometrically local Hamiltonian}
\label{sec:geo-local}
In this section, we consider a $d$-dimensional lattice $\Lambda \subseteq \mathbb{Z}^d$ and assume that $\ad_H = \sum_{\gamma\in [\Gamma]} \A_{\gamma}$ is geometrically local. 
For any $x\in \Lambda$ and a positive number $R$, define $B_R(x)$ as the $\ell_1$ ball centered on $x$ with radius $R$: \begin{align*}
    B_R(x) = \{ y\in \Lambda:d(x,y) \le R\},
\end{align*}
where \begin{align*}
    d(x,y) = \sum_{i=1}^d  |x_i-y_i|
\end{align*}
is the $\ell_1$ distance between $x = (x_1, \ldots, x_d)$ and $y= (y_1, \ldots, y_d)$. 
Assume that we can group the summand in the decomposition of $\ad_H$ and $\L_{\rm tot}$ as \begin{align*}
    &\ad_H = \sum_{x\in {\Lambda}} \A_x, \quad \A_x = \sum_{\gamma\in \mc I_x} \A_{\gamma} 
\end{align*} where $\{\mc I_x\}_{x\in \Lambda}$ is a partition of $[\Gamma]$, such that \begin{align*}
    &\supp(\A_{\gamma}) \subseteq B_R(x), \quad \forall \gamma\in \mc I_x
\end{align*} for some positive number $R$. We also assume that the observable of interest $O$ is supported on $B_{R_O}(x_O)$ for some positive $R_O$ and $x_O\in \Lambda$. Define \begin{align*}
    J_H = \max_{x\in \Lambda}\sum_{\gamma\in \mc I_x} \|\A_{\gamma}\|.
\end{align*} Define $\Lambda_d = 2^d/d!$ such that \begin{align*}
    |B_R(x)| = \Lambda_dR^d + O(R^{d-1})\le C_d R^d, \quad \forall x\in \Lambda,
\end{align*}
for some $C_d\le 2\Lambda_d$.
Then $\ad_H$ and $\L_{\rm tot}$ are $g_H$ and $g_L$-extensive with \begin{align*}
    g_H &\le \max_{x\in \Lambda} \sum_{y\in B_R(x)}\sum_{\gamma\in \mc I_y} \|\A_{\gamma}\|\le |B_R(x)|J_H  = O(C_d R^d J_H).
\end{align*}
For any $l \ge 0$, define\begin{align*}
    \ad_{H_{(l)}} = \sum_{y\in B_{R_O+l+R}(x_O)} \A_y,
\end{align*}
be the truncated Hamiltonian terms within a ball around the support of $O$. The difference between the time-evolved observables $O$ under $\ad_{H_{(l)}}$ and $ \ad_H$ can be bounded as follows. 
\begin{lemma}[{\cite[Lemma 2]{rao2025stability}}]
    For any $l\ge 0$, we have \begin{align*}
        \left\| e^{\i \ad_H t}O - e^{\i \ad_{H_{(l)}}t}O\right\| \le |\supp(O)|\|O\|\min\{e^{-l/R}(e^{eC_dR^d t}-1), 1\}.
    \end{align*}
\end{lemma}
We can choose \begin{align*}
    l = eC_d R^{d+1}T+ R\log(2|\supp(O)|/\eps),
\end{align*}
such that for any initial state $\rho_0$ \begin{align*}
    \left| \operatorname{Tr}[Oe^{-\i \ad_H T}\rho_0]- \operatorname{Tr}[Oe^{-\i \ad_{H_{(l)}} T}\rho_0]\right| \le \left\| e^{\i \ad_HT}O - e^{\i \ad_{H_{(l)}}T}O\right\| \le \eps/2.
\end{align*}
Therefore, we can simulate the truncated Hamiltonian $H_{(l)}$ on $B_{R_0+l+2R}(x_O)$,  and the noise terms are restricted to $B_{R_O+l+2R}(x_O)$. Define
\begin{align*}
     \L_{(l), \rm tot} =  \sum_{v\in B_{R_O+l+2R}(x_O)} \L_v,
\end{align*}
which includes all possible noise terms. 
\xz{TBD}
\sz{The intended next step is to combine the locality truncation above with \thm{extrapolation-k-local}. Let
\begin{align*}
    N_{\mathrm{eff}} = \big|B_{R_O+l+2R}(x_O)\big| = O\!\left(C_d(R_O+eC_dR^{d+1}T+R\log(|\supp(O)|/\eps)+R)^d\right)
\end{align*}
be the effective number of sites in the truncated light cone. Applying \thm{extrapolation-k-local} to the truncated Hamiltonian $H_{(l)}$ and the restricted noise generator $\L_{(l),\rm tot}$ gives an estimator for $\operatorname{Tr}[Oe^{-\i\ad_{H_{(l)}}T}\rho_0]$ with error at most $\eps\|O\|/2$. Together with the truncation error above, this yields an $\eps\|O\|$-accurate estimator for the original local observable. The resulting complexity should be obtained from \thm{extrapolation-k-local} by replacing the global system size $N$ with $N_{\mathrm{eff}}$, so the dependence on the total lattice size disappears and is replaced by the volume of the effective light cone around $\supp(O)$.}

\section{Experimental Demonstration}

The Hamiltonian used for the experimental demonstration is a one-dimensional transverse-field Ising model,
\begin{equation*}
    H = -J \sum_{\langle i,j\rangle} Z_i Z_j - h \sum_{i=1}^{n} X_i,
\end{equation*}
with periodic boundary conditions on a 9-qubit ring. The system is evolved for a total time $T$ using a second-order symmetric (i.e., $p=1$) Trotter decomposition, \sz{numerics to align with the topological structure.}
\begin{equation*}
    \mathcal{P}(s) = e^{-i h s X^{\otimes n}/2} e^{-i J s \sum Z_i Z_j} e^{-i h s X^{\otimes n}/2},
\end{equation*}
where $s = T/r$ is the Trotter step size and $r$ is the number of Trotter steps.

The observable measured throughout the hardware experiments is the average longitudinal magnetization,

\begin{equation*}
O=\frac{1}{n} \sum_{i=1}^n Z_i,
\end{equation*}
which can be estimated directly from computational-basis measurements.

\subsection{Real quantum hardware experiments}
We conduct physical experiments on the 13-qubit Tianji S2 superconducting quantum chip. We adopt the sparse Pauli-Lindblad model in~\cite{berg_probabilistic_2023} to characterize the physical noise, which assumes that each layer of the quantum circuit is followed by
\begin{equation}\label{eq:lind}
\mathcal{L}(\rho)=\sum_{k \in \mathcal{K}} \lambda_k\left(P_k \rho P_k-\rho\right).
\end{equation}
In particular, we assume that $|\mathcal{K}| \ll 4^n-1$ is a set that is only of polynomial size in the number of qubits. The final noise channel is given by the exponential of \eqref{eq:lind} as:
\begin{equation*}
\Lambda(\rho)=\prod_k\left(w_k \cdot+\left(1-w_k\right) P_k \cdot P_k^{\dagger}\right) \rho,
\end{equation*}
where $w_k=\left(1+e^{-2 \lambda_k}\right) / 2$. \sz{Quantum circuit schematic figure. $r_x=-h$, $r_{zz}=-J$.}

We first learn the coefficients $\lambda_k$ based on an adaptive/improved scheme
\begin{figure}[!htbp]
    \centering
    \includegraphics[width=0.92\linewidth]{figures/archive/trotter_ising_hardware_implementation.pdf}
    \caption{
    Hardware implementation of Trotterized Ising dynamics.
    (a) A second-order symmetric Trotter step alternates transverse-field rotations and Ising interactions.
    (b) The interaction layer is decomposed into three parallel $R_{ZZ}$ patterns.
    (c) A connected 9-qubit subgraph of the 13-qubit device is used; colored edges show the scheduled interaction patterns.
    (d) Logical $R_X$ and $R_{ZZ}$ operations are compiled into the native $H$, $R_Z$, and controlled-$Z$ gates.
    }
    \label{fig:hardware-implementation}
\end{figure}

\begin{figure}[!htbp]
    \centering
    \begin{subfigure}{0.48\linewidth}
        \centering
        \includegraphics[width=\linewidth]{figures/archive/hardware_joint_extrapolation_minus1.pdf}
        \caption{}
    \end{subfigure}
    \begin{subfigure}{0.48\linewidth}
        \centering
        \includegraphics[width=\linewidth]{figures/archive/hardware_joint_extrapolation_minus2.pdf}
        \caption{}
    \end{subfigure}
    \caption{
    Combined hardware extrapolation of noise and Trotter errors.
    Numerical simulation (blue) and superconducting-hardware data (orange) are plotted versus $x=1/N_{\mathrm{Trotter}}$.
    The fitted curves extrapolate toward $x=0$, and the red triangle marks the exact reference.
    Panels show two representative evolution settings: (a) $r_xT=r_{zz}T=-1$ and (b) $r_xT=r_{zz}T=-2$.
    }
    \label{fig:hardware-joint-extrapolation}
\end{figure}


\subsection{Classical tensor network simulation}
\sz{In addition to the hardware experiments, which are limited to a relatively small number of qubits by device fidelity, we use classical tensor-network simulations to study larger one-dimensional Ising instances and to compare different extrapolation strategies in a controlled setting.}
In this section, the results and plots are obtained by numerical simulations on a local machine with an Apple M2 CPU (8 cores) and 16 GB unified memory via Python 3.11.12. 
We use the `quimb' library~\cite{Gray2018} to represent (noisy) quantum circuits using tensor networks.
The target observable is again the average longitudinal magnetization. We compare three mitigation strategies: the proposed joint one-dimensional extrapolation, a noise-free Trotter extrapolation that removes only the algorithmic Trotter error, and a two-dimensional extrapolation baseline that treats the Trotter step size and the physical noise strength as independent extrapolation variables.
The goal of this comparison is to test whether coupling the physical noise strength to the Trotter step size can retain the accuracy of the two-dimensional strategy while using a one-dimensional extrapolation procedure.
For the two-dimensional baseline, we use a tensor-product Richardson extrapolation. The Trotter direction uses the four step counts $r\in\{1,2,3,7\}$, equivalently step sizes $1/r$, and the noise direction uses the four positive noise-amplification factors $\{49,9,4,1\}$, \sz{in practice, explain why we instead use squared instead of cubic} extrapolated to zero noise. Thus, the Richardson order in each direction is the number of data points in that direction, and the two-dimensional baseline uses $4\times 4=16$ noisy data points. The noiseless curve labeled ``Noise 0'' is included as a reference and is not counted as a noisy ZNE data point.
\textcolor{red}{Xiaoyang: since the tensor network is limited to small depth $r$, we can instead polt log error - $T$, left: before extrapolation maybe use $r_m=7$, right after extrapolation.}

\begin{figure}[!htbp]
    \centering
    \includegraphics[width=0.8\linewidth]{figures/archive/richardson_error_vs_r_N_20_seed_43.pdf}
    \caption{Average magnetization error as a function of the number of Trotter steps $r$ for a $1$D transverse-field Ising model with $N=20$ qubits. The curves compare the raw noisy estimates at several noise-amplification levels with the estimates obtained by the joint one-dimensional extrapolation, the noise-free Trotter extrapolation, and the tensor-product two-dimensional Richardson baseline. Results are obtained by classical tensor-network simulation with \textup{\texttt{quimb}} at seed 43.}
    \label{fig:richardson-error-vs-r}
\end{figure}

\begin{figure}[!htbp]
    \centering
    \includegraphics[width=0.8\linewidth]{figures/archive/richardson_error_vs_inv_r_N_20_seed_43.pdf}
    \caption{The same average magnetization error data plotted against the Trotter step size $1/r$. This view highlights the extrapolation behavior as the step size approaches zero. The inset compares the final extrapolated estimates from the joint one-dimensional method, the noise-free Trotter extrapolation, and the two-dimensional baseline. Results are obtained by classical tensor-network simulation with \textup{\texttt{quimb}} at seed 43.}
    \label{fig:richardson-error-vs-inv-r}
\end{figure}

It can be seen from \fig{richardson-error-vs-r} and \fig{richardson-error-vs-inv-r} that the proposed one-dimensional joint extrapolation substantially reduces the bias of the noisy estimates across the simulated Trotter step sizes. In the displayed instance, the final extrapolated value obtained by the one-dimensional method achieves comparable accuracy to the two-dimensional Richardson baseline while using a one-dimensional extrapolation procedure. This supports the central intuition of the protocol: once the physical noise strength is coupled to the Trotter step size at the appropriate order, a single extrapolation variable can remove the leading physical-noise and Trotter contributions simultaneously.

\section{Conclusion and discussion}

We have developed a commutator-sensitive analysis of Trotter extrapolation in the presence of physical noise. Our analysis builds on the finite-order BCH truncation bound from Ref.~\cite{wang2026lindbladiansimulationcommutatorbounds}, which is controlled by doubly right-nested commutators and avoids requiring convergence of the full BCH series. In the present work, we use this tool to analyze noisy product formulas and to extend the commutator-scaling guarantees of Trotter error mitigation to a joint mitigation protocol for algorithmic and physical errors.

The main contribution of this paper is a one-dimensional joint extrapolation protocol for simultaneously mitigating algorithmic Trotter error and physical noise. By coupling the noise strength to the Trotter step size as $\lambda(s)=c(sT)^{p+1}$ for a $p$-th order product formula, the leading Trotter and noise contributions enter at the same order in the extrapolation variable. Richardson extrapolation can therefore remove both contributions using a single set of extrapolation nodes. For $k$-local and geometrically local Hamiltonians under sufficiently mild local noise, our bounds show that physical noise can be mitigated together with Trotter error with only a constant asymptotic overhead compared with the corresponding noiseless Trotter extrapolation protocol.

Several assumptions are important for interpreting these results. The protocol requires access to a tunable family of effective noise strengths, at least in the range needed by the extrapolation nodes. In practice, this tuning may be implemented by noise amplification techniques such as unitary folding, or by model-dependent procedures based on a learned sparse Pauli--Lindblad noise model. The achievable range and accuracy of this tuning depend on the hardware platform and on the quality of noise characterization. Our theoretical analysis also focuses on local Lindbladian noise and on expectation-value estimation; extending the same framework to more general non-Markovian noise models, correlated readout errors, or adaptive measurement strategies remains an interesting direction.

The numerical and hardware demonstrations support the basic mechanism of the protocol, but they are not intended to exhaust the possible performance regimes. Future work could optimize the choice of extrapolation nodes and the coupling constant $c$, quantify constant-factor overheads for concrete devices, and test the protocol on larger Hamiltonian simulation tasks where Trotter error and physical noise are both visible at experimentally accessible circuit depths.

\bibliographystyle{qubit}
\bibliography{ref-codex}
\appendix

\end{document}
